% PaceNotes mindmap, section 5: detailed research record (user request 2026-09-27 KST: the mindmap keeps
% "we tried this -> this happened -> so we decided this" in detail; the main paper is polished separately).
% Thread files 5_rec_a..f.tex and 5_record_overview.tex are generated by D:/tools/scratch_qdd/mmrec/render.py.
\section{연구 기록 --- 해 봤다, 이렇게 됐다, 그래서 이렇게 하기로 했다}
\label{sec:record}

\noindent 이 절은 연구 줄기마다 시간순으로 적은 상세 기록이다(2026-09-27 KST 작성). 앞 1--4절은 모듈별 생각의 흐름을 그대로 두고, 여기서는 한 줄기의 결정이 어떻게 이어졌는지를 한곳에서 따라가게 한다.

\paragraph{읽는 법}
\begin{itemize}
  \item 항목마다 \textbf{사용자}(빨강, \texttt{docs/user-log.md} 원문 그대로) $\rightarrow$ \textbf{해 봄}(실험·설정·사전 등록 커밋) $\rightarrow$ \textbf{결과}(결과 문서의 숫자 그대로, 커밋 해시) $\rightarrow$ \textbf{그래서(결정)}(정본 절, 채택·불채택) $\rightarrow$ \textbf{다음} 순서다.
  \item \tentative{주황}은 아직 진행 중이거나 결과가 없는 것이다. 숫자를 지어내지 않는다 --- 결과 문서가 없는 실험은 ``진행 중''으로만 적는다.
  \item 시각은 모두 KST(UTC+9)다. 저장소의 원래 기록(정본·user-log·결과 문서)은 UTC이고, 커밋 시각은 git 기록(KST) 그대로다. ``ul N''은 user-log 번호, ``§N''은 정본 \texttt{docs/design/00-interfaces.md}의 절이다.
  \item 통제자가 전한 사용자 말도 user-log에 번호를 붙여 올린 뒤(ul 130--133) 빨강으로 적는다. 빨강은 모두 user-log 원문과 글자가 같아야 한다(\texttt{tools/intent\_check.py}).
\end{itemize}

% Generated by D:/tools/scratch_qdd/mmrec/render.py from data_a.py - edit the data, not this file.
\subsection{A. Astra 활용 --- 무엇을 맡기고, 한 번의 답은 얼마나 맞고, 돈은 얼마를 쓰나}\label{rec:A}
\noindent\textbf{한 줄:} Astra만 할 수 있는 일로 역할을 좁히고(§82), 탐침으로 흐름 상수를 정하고(§86), 한 호출 정확도를 재어 v2 프롬프트를 채택했지만 7--10 cm 어긋남을 못 알아채는 병목(0/13)이 남았다. Astra에 맞춘 단독 조종은 7/7 성공. 유료 호출은 이제 사용자 명시 승인 뒤에만 한다.

읽는 법: 각 항목은 사용자 말(빨강) $\rightarrow$ 해 봄 $\rightarrow$ 결과 $\rightarrow$ 그래서(결정) $\rightarrow$ 다음 순서다. 커밋 해시와 시각(KST)은 저장소 기록 그대로이고, 숫자는 결과 문서 값 그대로다.

\paragraph{A1. Astra를 쓰는 이유가 있어야 한다 $\rightarrow$ Astra 전용 일 J1--J6}
{\small 09-25 23:15 -- 23:38 KST}
\begin{itemize}
  \item \textbf{사용자} (ul 74, 09-25 23:15 KST): \intent{그 아스트라라는 존재를 상위 오케스트라로 쓰는 이유가 존재해야만해 얘가 자기 혼자서도 로봇조작을 할 수 있는놈이기때문에 아스트라를 너무 별일 아닌데도 활용을 해버리면, 그냥 더 작은 vlm쓰는게 맞음 최대한 엄청 최대한 능력치를 다써야해 아스트라를 쓸 수 밖에 없는 이유를 만들 수 있는 장점발휘를 할 수 있게 세팅을 해야함 좀 찾아서 어케 활용할지를 더 구체화해봐봐}
  \item \textbf{해 봄:} 조사 \texttt{docs/\allowbreak{}research/\allowbreak{}astra\_\allowbreak{}role\_\allowbreak{}2026-\allowbreak{}09-\allowbreak{}25.\allowbreak{}md}(925b0b0, 23:34): Astra만 할 수 있는 일, 작은 모델로 내릴 일, 호출 사다리, 필요성 실험 E-Astra-necessity 초안.
  \item \textbf{결과:} 강등 근거(정본 §82): V1h 세계 술어 0.949·0.44 ms 대 Astra 첫 토큰 약 3 s; FORTRESS(CoRL 2025) 즉석 GPT-4o 판정 0.64 대 사전 계산 경로 0.90.
  \item \textbf{사용자} (ul 76, 09-25 23:38 KST (질문: 하트비트·단계 경계 확인을 작은 모델로 내리고 Astra는 J1--J6에만 쓸지) 선택): \intent{바꾸기 (Recommended)}
  \item \textbf{사용자} (ul 76, 09-25 23:38 KST): \intent{논문에 로우 하이를 모두 가져가는 쪽이 맞겠네 성능비교두 그렇고}
  \item \textbf{사용자} (ul 76, 09-25 23:38 KST): \intent{지금 10만원 정도는 비용 한도 책정이 가능해}
  \item \textbf{그래서(결정):} 정본 §82(a1bb055, 23:38): Astra는 J1 과제 컴파일·J2 실패 진단과 복구안·J3 검사 술어 합성·J4 새 물체 이름·J5 저빈도 진행 감사·J6 경험 증류에만 쓴다. 5 s 하트비트는 J5로 대체, effort는 low·high를 모두 논문에 보고, 유료 한도 약 10만 원(4eb988d).
  \item \textbf{다음:} E-Astra-necessity(같은 자리에 로컬 VLM) 사전 등록 초안(2eea17b, 09-26 13:42)은 유료 실행 금지 상태로 남아 있다(A10).
\end{itemize}

\paragraph{A2. Astra가 움직임에도 개입 $\rightarrow$ 흐름 교정 설계, 호출은 한 번에 하나}
{\small 09-26 00:16 -- 02:18 KST (+ 14:14)}
\begin{itemize}
  \item \textbf{사용자} (ul 78, 09-26 00:16 KST): \intent{나는 그냥 아스트라 자체가 이 움직임에도 개입을 했으면 좋겠어 근데 딜레이가 문제니까 이걸 해결하면서 도 개입을 하는 방법을 찾는거야 최소화 하는 것도 방법이겠지 VLA가 한거를 감사하는건 아스트라의 능력중 극히 일부일 뿐인거구}
  \item \textbf{사용자} (ul 78, 00:18 KST (A: M4 느린 투표자 / B: 물체 기준 경로 / C: 조건부 프로그램)): \intent{A,B 더하는 방식이 조금 더 좋은 것 같긴하다}
  \item \textbf{사용자} (ul 83, 01:17 KST): \intent{로우로하자 그럼 이 로우와 하위 vla가 rtc로서 서로서로 보완하면 좋을듯 아스트라는 반응성이 느려서 그러지 vla는 방향성이 빠르지만 자기가 뭘하는지에대한 성능이 조금 떨어져서 그렇지}
  \item \textbf{사용자} (ul 83, 01:21 KST): \intent{이게 뭐랄까 엄청 촘촘하게 이루어질 수 있어야하는데 api를 직렬이 아니라 1초단위로 3개 불러오면 1초1초가 연속적인 방법이 될 수는 있잖어}
  \item \textbf{사용자} (ul 83, 01:33 KST): \intent{목표를 공간에 표현하는 것 조차 안하고 그냥 일반적 아스트라 조종법을 따르겠다는 얘기}
  \item \textbf{사용자} (ul 83, 01:33 KST): \intent{네, 일반 조종법만}
  \item \textbf{사용자} (ul 83, 01:40 KST): \intent{아스트라는 3개의 캠을 다봐야해 예를들어 헤드캠에서 봤을때는 원시때문에 물체를 잡을 때, 잡았다고 오탐하지 않게 우론목에는 안보이는 그런게 있는거지 그걸 해결할 수 이썽야하는거지}
  \item \textbf{해 봄:} 결합 설계 초안(2cdef56, 01:26): 손끝 목표 파이프라인, 1 Hz 계단식 Astra low, 두 층 M4, RTC 이음. 개정(ac165ab, 01:33): 공간 목표 표현 없이 일반 조종(continue/edit/stop + 평가). 설계 §11 부드러운 반영(00a0aad), §12 카메라 3대와 손목 근거(f71f407).
  \item \textbf{그래서(결정):} 정본 §82 보충 2(a8259df, 01:17): 실행 중 Astra effort = low, Astra와 VLA는 RTC식으로 서로 보완. high는 비교 실험·오프라인 일에서만.
  \item \textbf{사용자} (ul 86, 02:18 KST): \intent{그냥 아스트라 부르는거 그냥 1번씩만 하는걸로하자}
  \item \textbf{그래서(결정):} 정본 §84 보충 2(190af17, 02:18): 흐름 호출은 동시 요청 1개 --- 답(또는 시간 초과)이 오면 그때의 최신 관측으로 다음 요청. 1 s 계단식·동시 $\leq$ 3을 대체. `동시 1개'는 Claude 해석이며 사용자에게 알렸다(에피소드당 한 번이라는 뜻이면 정정).
  \item \textbf{사용자} (ul 96, 09-26 14:14 KST): \intent{주기로 아스트라 부르는건 끄는거?}
  \item \textbf{사용자} (ul 96, 09-26 14:14 KST): \intent{그 주기로 부르는건 켜두고.}
  \item \textbf{그래서(결정):} 흐름 호출은 켠 채 유지(17f7513). 새로움 게이트(E-NOV0)는 선택지 시험일 뿐이고 호출 방식 변경은 사용자 결정.
\end{itemize}

\paragraph{A3. E-Astra-motion 탐침 설계 --- 돈을 쓰기 전에 유익한가부터}
{\small 09-26 00:28 -- 03:27 KST}
\begin{itemize}
  \item \textbf{사용자} (ul 79, 00:28 KST): \intent{탐침 E-Astra-motion: 3D 경로 그리기 대 보고 조종하기 이거 한번 해보자용}
  \item \textbf{사용자} (ul 79, 00:31--00:37 KST): \intent{뎁스를가지고 한다고 하면 될 것 같아 뎁스가 있잖아 ai워커가}
  \item \textbf{사용자} (ul 79, 00:31--00:37 KST): \intent{아니다 뎁스를 쓰면 너무 범용성이 떨어질 것 같은 생각이드네 이미지 상만으로 어떻게 3D를 방향을 제대로 계획을 시킬 수 있을까...?}
  \item \textbf{사용자} (ul 79, 00:31--00:37 KST): \intent{그냥 뎁스라기보다는 그 가우시안 스플레팅인가? 아니면 포인트 클라우딩 포인트 클라우딩으이었던거 같은데 스테레오캠으로 하는 그거 그거를 어떻게 잘 해서 안되는건가? 공간 좌표 만들고 거기서 하는거.?}
  \item \textbf{해 봄:} 사전 등록 a62164d(02:39, 문서 시각 17:38:25Z, 유료 호출 전): 직렬 흐름, 요청 방식 F0/F1, 잡기 판정(카메라 구성), 부드러운 반영.
  \item \textbf{사용자} (ul 87, 02:41--02:45 KST): \intent{그 돈쓸떼는 정말 이게 유익한 실험이 되는지도 중요해}
  \item \textbf{사용자} (ul 87, 02:41--02:45 KST): \intent{고 나에게 승인은 안받아도되지만 자체적으로 저걸 해보라는거지}
  \item \textbf{사용자} (ul 87, 02:41--02:45 KST): \intent{그리고 만약 중간에 잘못되면 다시 새롭게 설계헤ㅐ서 낭비 없게 해야해}
  \item \textbf{사용자} (ul 87, 02:41--02:45 KST): \intent{매번 자체검사를 해야해 알겠지?}
  \item \textbf{결과:} 유료 일시 정지(5320dd2, 02:41). 옛 설계 G1은 13호출 204원(호출당 15.7원, 지연 p50 4.2 s). 무료 Qwen 사전 실행에서 카메라 3대일 때 `잡았다' 과잉 주장 12/14 $\rightarrow$ 원인은 `쥠'을 정의하지 않은 프롬프트, 정의 뒤 1/7(정본 §86).
  \item \textbf{그래서(결정):} 재범위(075c505, 02:43): 무료 Qwen 먼저, 유료는 G1(약 40장·4조건) + S2(F0·F1 각 3편) + G2(low 오답에만 high $\leq$ 8회), S1에서 Astra 뺌, 누적 하드 정지 15,000원; 관문마다 결함이면 멈추고 재설계. 자체 검사 규칙을 CLAUDE.md로(ed478e4, 02:45).
  \item \textbf{참고:} 등록 수정 1·2 커밋(35029e5, 03:27)은 유료 호출 재개보다 늦게 커밋되었다 --- 이탈로 기록(R7 25회차 E-D2).
\end{itemize}

\paragraph{A4. 탐침 결과 $\rightarrow$ 결합 상수 확정(§86)}
{\small 09-26 04:09 KST}
\begin{itemize}
  \item \textbf{결과:} 요청 방식: 뒤집힘율 0.594 $\rightarrow$ 0.381($-$36 \%, 등록 규칙 50 \% 감소 미달). 직렬 흐름(카메라 3대) 지연 p50 약 9--10 s·p95 12--16 s, 답 나이 p50 6.7--9.8 s --- 6 s 문턱이면 답의 70--100 \%가 버려진다.
  \item \textbf{결과:} 잡기 판정 40장(쥠 25/안 쥠 15): 머리만이면 uncertain 39/40, 손목 포함이면 오탐 0--1/15이지만 실제 쥠 13/25를 놓침(한 손가락 테두리 걸기 6/6 전부), high 표적 8장 중 2장만 고침; 덧그림 효과 없음(짝 3 대 3); `불확실 $\rightarrow$ continue' 규칙 발동 0회.
  \item \textbf{결과:} 부드러운 반영(개루프 재생): 저크 RMS 중앙 0.83 대 즉시 2.36 m/s$^3$, 최대 속도 0.036 대 0.095 m/s. Astra 단독($\leq$ 5 cm 수정, VLA 없음)은 물체 도달 0/6, Qwen 동기 0/20. 비용: 흐름 49.7--58.0원/호출, 파지 질문 16.8원, high 표적 27.9원, 로봇 1분당 약 350--500원, 캐시 적중 0, 탐침 총 6,933.4원(하드 정지 15,000원).
  \item \textbf{그래서(결정):} 정본 §86(3500f54): 요청 방식 = F0; 늦은 답 폐기 6 $\rightarrow$ 15 s(보충: 시간 초과 20 s); 파지 확정은 Astra 영상이 아니라 고유감각 T1·V1h(Astra 주장은 참고), 손목 영상은 Astra 입력에 유지; 확신 필드로 거르지 않고 두 답 합의·편향 속도 상한으로; 부드러운 반영 유지; 모든 판정 질문 프롬프트에 판정 대상의 조작적 정의. \texttt{CoupleParams}: F0·15 s·20 s·\texttt{latency\_\allowbreak{}init\_\allowbreak{}s} 9.3 s·\texttt{est\_\allowbreak{}text\_\allowbreak{}tokens} 1601(§86 보충 2).
  \item \textbf{다음:} Astra 단독 0/6은 뒤에 사용자가 `실험을 Astra에 맞게 안 한 탓'이라고 지적 $\rightarrow$ A8.
\end{itemize}

\paragraph{A5. 프롬프트 검진}
{\small 09-26 14:14 -- 15:51 KST}
\begin{itemize}
  \item \textbf{사용자} (ul 96, 14:14 KST): \intent{프롬프트 엔지니어링 등으로 문제가 생겼는지 검진하는게 필요해보이는데}
  \item \textbf{해 봄:} 모든 프롬프트(Astra 흐름·J1--J6·탐침 판정 질문·로컬 VLM 어댑터·VLA 직렬화 문장)를 정적 감사하고, 무료 Qwen3-VL 8B·4B로 문구·보기 순서·카메라 순서·정의를 흔들고, 유료 Astra 표본(사전 등록 5b62187, 14:39, 상한 3,000원·80 \% 정지).
  \item \textbf{결과:} 1b68a6a(15:16): \texttt{astra-\allowbreak{}couple@v1}에 판정 정의·축 문구가 없고, 그리지 않은 화살표를 범례가 설명(Qwen 8B 인용 G 27/33), 확정 화살표 = 결정 토큰 중심값. Qwen 문구 뒤집힘 8B 6--17/40·4B 최대 28/40(같은 문구 반복 0/40), 방향 질문은 우연 수준. Astra 잡기 문구 뒤집힘 3/30(강건), 축 문구 18/30 대 17/30(가를 수 없음, x 부호 오류 = 원근). VLA \texttt{mag\_\allowbreak{}coarse}·\texttt{dir\_\allowbreak{}xy} 라벨 뜻 $\neq$ 질문 문구(청크 일치 0.375·0.626) $\rightarrow$ E-SR0·E-SR1b의 부분 원인. 탐침 Astra 답의 도착 시점 방향 F0 25/40·F1 20/24. 2,093.9원, 약 0.8 GPU-h.
  \item \textbf{그래서(결정):} 운영 프롬프트는 그대로 두고(판본 기록 누락 1건만 고침) 발견은 새 판에: Astra 쪽은 \texttt{astra-\allowbreak{}couple@v2}(A6), VLA 쪽은 PH-A1(b931f75, 15:51) --- 라벨은 그대로, 융합 VLA 문구를 라벨 정의대로. 로컬 VLM 이미지 상한 2 $\rightarrow$ 3(§59 보충, a913020).
\end{itemize}

\paragraph{A6. Astra 한 번의 명령은 정확해야 한다 $\rightarrow$ E-ACC 1단계, v2 채택}
{\small 09-26 14:24 -- 19:41 KST}
\begin{itemize}
  \item \textbf{사용자} (ul 97, 14:24 KST): \intent{아스트라가 다음 답변을 할때까지 아무것도 못하는 거니까. 아스트라가 한번 명령을 내릴때 좀 정확히 해야함}
  \item \textbf{사용자} (ul 104, 15:24 KST): \intent{아스트라를 최대 성능이 나게 만들어놓는것이 광장히 중요함}
  \item \textbf{사용자} (ul 104, 15:24 KST): \intent{vla는 아스트라가 옳다는 전제로 움직이기때문}
  \item \textbf{그래서(결정):} 정본 §94(64b0a45, 15:24): Astra 호출 단위 정확도 최대화가 E-Couple의 선행 조건, E-ACC를 먼저.
  \item \textbf{해 봄:} 사전 등록 e4f8891(15:52): Isaac 벤치(R2 오라클 계획기 $\times$ 0.5 속도를 대역 정책으로), 스냅숏 종류 on·off\_a(7--10 cm 엉뚱한 곳)·off\_b(엉뚱한 물체)·off\_c(엉뚱한 놓을 곳), 팔 v1·v2·v2cp·v2\_ax, 채점 = 도착 시각(스냅숏 + 9.3 s) 기준. 벤치 42편. 무료 Qwen 선별 G-fmt 실패 $\rightarrow$ 변경 2(eb42189)로 형식만 고침.
  \item \textbf{결과:} 사고: P1의 31번째 호출부터 OpenAI 크레딧 소진(\texttt{insufficient\_\allowbreak{}quota}) --- 빈 답 37개, 과금 1,352.7원(on 30호출만). 클라이언트가 스트림 \texttt{error}·\texttt{response.\allowbreak{}failed}를 오류로 올리지 않았다(함정 P108, cfaf140 16:45).
  \item \textbf{그래서(결정):} 변경 3·4(a09c5ed, 19:05): 새 장부(과금 이월), quota 즉시 정지·빈 답 3연속 정지, off 20장 $\times$ v1·v2·v2\_ax로 재개, v2cp는 뺌.
  \item \textbf{결과:} 2a7b9ae(19:15): 무효 0/90. 수정(edit)을 냈을 때 방향은 13/13 맞음(도착 각 중앙 5--9$^\circ$). 그러나 off\_a·off\_c는 세 팔 모두 0/13(`progressing / aligned'). off\_b는 v1 1/7·v2 5/7·v2\_ax 7/7. R1 v2 $-$ v1: M1 +0.20 [+0.05, +0.40] $\rightarrow$ v2 채택. R2$'$ 축 덧그림: +0.10 [0.00, +0.25] $\rightarrow$ 미정, 기본 끔. 구간 now 정답 v2 26/30. 원/호출 v1 36.4·v2 49.1·v2\_ax 51.3.
  \item \textbf{그래서(결정):} \texttt{astra-\allowbreak{}couple@v2}(8c9064a, 19:41, 결합 Task 18) = E-ACC 채택 문장을 바이트 그대로(판본 83fa03a5de19): 축 줄, 정의(쥠 = 두 패드가 닫혀 물체에 붙음), 실제로 그린 요소만 범례, \texttt{segment\{now, do, next\}}, \texttt{valid\_\allowbreak{}until}, \texttt{since\_\allowbreak{}last\_\allowbreak{}request}. AxisGuide·카메라 자세 줄은 기본 끔.
\end{itemize}

\paragraph{A7. 알아채기 병목 --- E-ACC 2단계와 B$'$}
{\small 09-26 19:18 -- 19:34 KST}
\begin{itemize}
  \item \textbf{해 봄:} 변경 5(0dd95ad, 19:18): off\_a·off\_c 13장에서 (A) effort medium, (B) 목표 확인 문장(on 5장으로 괜한 수정도 확인).
  \item \textbf{결과:} 39ae001(19:28): medium 1/13, 지연 p95 17.6 s $\rightarrow$ 권고 안 함(low 유지). 목표 확인 7/13(+0.54 [+0.23, +0.77]; off\_a 6/7·off\_c 1/6)이지만 계획대로 상태 5장 중 괜한 수정 2장 $\rightarrow$ S2B 불채택. 괜한 수정의 원인은 `도착 예상 손끝' 외삽(0.5 s 속도 $\times$ 9.3 s)이 6--9 cm 튄 것(P113).
  \item \textbf{해 봄:} 변경 6(47e6ff6, 19:29): B$'$ --- 목표 확인의 기준점을 `지금 손끝 + 확정 화살표'로.
  \item \textbf{결과:} f54cc88(19:34): 괜한 수정은 사라졌지만(on 5/5) 알아채기 3/13(+0.23 [0.0, +0.46]), off\_c 0/6 $\rightarrow$ S2B$'$ 불채택. E-ACC 합계 6,892.2원(상한 8,000원).
  \item \textbf{그래서(결정):} 두 수단 모두 v2에 넣지 않는다. 대신 결합의 예상 상태 기본을 외삽이 아니라 `지금 손끝 + 실행 청크 변위'로(§91 구현, Task 19, 20d735c 20:38). 남은 병목: Astra는 7--10 cm 어긋난 목표·놓을 곳을 알아채지 못한다(0/13) --- VLA는 Astra가 옳다고 보고 움직이므로(§94) 결합의 가장 큰 오류 전파 경로로 기록(§90 구현 기록 3).
  \item \textbf{다음:} 보고만 한 후보: 외삽을 고친 뒤 B 다시 재기, 로봇 좌표 평면도(APC). 둘 다 새 등록과 상한 증액 검토가 필요해 실행하지 않았다.
\end{itemize}

\paragraph{A8. Astra 단독도 꽤 할 것 $\rightarrow$ Astra에 맞춘 조건에서 4/4}
{\small 09-26 15:18 -- 16:38 KST}
\begin{itemize}
  \item \textbf{사용자} (ul 103, 15:18 KST): \intent{아스트라 혼자조종도 꽤 할거야. 나는 vla와 아스트라의 하이브리드가 그리 옳다고는 생각안해}
  \item \textbf{사용자} (ul 103, 15:19 KST): \intent{3이 맞음 근데 아스트라가 0이 나온건 너가 최적으로 아스트라에 맞게 실험을 안했다는 얘기기도해}
  \item \textbf{그래서(결정):} 정본 §93(ed06d8a, 15:19): 하이브리드 유지. 탐침 0/6은 Astra에 불리한 조건(v1 프롬프트 결함, 수정 $\leq$ 5 cm, 90 s·30호출, 기다리지 않는 흐름)의 값이라 Astra 실력으로 인용하지 않는다. Astra 단독 팔은 Astra에 최적으로 다시 설계. 함정: 불리한 조건의 기준선은 비교가 아니다(P101).
  \item \textbf{해 봄:} 사전 등록 50336ce(16:07, 변경 1 포함, 첫 유료 호출 전): 동기(로봇이 답을 기다림), \texttt{astra-\allowbreak{}solo@v2}(정의·축·카메라 자세·그린 요소만 범례·측정 이력), 명령 eef 절대 목표 + edit $\leq$ 10 cm + gripper + stop, 최소 저크 실행기(VLA 없음), 머리 영상에 탁자면 로봇 좌표 격자(5 cm) + TCP 낙하선, 편당 40호출·180 s. 무료 단계: 참값 모델 6/6, Qwen3-VL-8B 0/5(인식 한계 --- 배선 결함 2개는 유료 전에 고침).
  \item \textbf{결과:} cd80896(16:38): Astra low 4/4(mug\_tray DEV standard 0·1, dr 0·1). 편당 9--17호출(평균 11.0), 성공까지 움직임 중앙 16.35 s, 벽시계 74--215 s, 성공 1건당 591.5원, 무효 0/44. dr 0은 첫 닫기가 49 mm 빗나간 것을 패드 간격으로 알아채고 다시 잡아 성공. 접근 목표 xy 오차(편 중앙의 중앙) 14.1 mm. effort 짝 12개 중 동률 아닌 짝 5 < 8 $\rightarrow$ undecided(medium 4승). 3,080.9원.
  \item \textbf{그래서(결정):} D1 ADOPT $\rightarrow$ 전체 팔은 \texttt{astra-\allowbreak{}solo@v2} 그대로; D2 $\rightarrow$ low 기본 + medium은 판정 밖 부분집합. 해석: E-Couple 사전 실행의 VLA 단독 0/12와 함께 보면 `Astra 단독(동기)'이 강한 기준선일 수 있다(§92 P3와 함께 봐야 함).
\end{itemize}

\paragraph{A9. Astra 단독 전체 팔 --- 크레딧 소진, 비용 지적, 멈춤}
{\small 09-26 18:42 -- 19:21 KST}
\begin{itemize}
  \item \textbf{해 봄:} 등록 변경 2(d0d2ae7, 18:42): AS 팔을 §92 P3 기준선으로 --- 40편(DEV 0--19 $\times$ standard/dr), 하드 정지 25,000원; 러너 가드(\texttt{insufficient\_\allowbreak{}quota} 즉시 치명, 빈/오류 답 3연속 정지). OpenAI 크레딧 소진으로 처음에는 보류(HOLD).
  \item \textbf{결과:} 19:00 KST(10:00 UTC)에 시작해 3편 모두 성공(9·9·9호출, 460·456·443원), 장부 1,359.0원에서 메인이 멈춤.
  \item \textbf{사용자 말(user-log 번호 없음)} (메인 전달, 등록 변경 3에 인용): ``한 번에 23,000원 쓰는 건좀 아닌 것 같은데 좀 아껴서 하지 제대로 된 결과 나오게 아껴서 하지''
  \item \textbf{그래서(결정):} 등록 변경 3(11cbc39, 19:13): 단계식 재설계 --- 1단계 = DEV 0--5 $\times$ standard/dr 12편(이미 3편), 1단계 하드 정지 8,000원; 러너 조기 종료(20호출 또는 움직임 60 s, 프롬프트 바이트는 그대로)로 실패 편 최대 약 2,100 $\rightarrow$ 1,050원; 순차 규칙($\geq$ 10/12 강함·$\leq$ 2/12 약함에서 멈춤). 비용 분해: 입력 약 63 \%, 캐시 적중 0/80.
  \item \textbf{진행 중·대기:} \tentative{1단계 나머지 9편은 결과 문서가 없다(실행 기록 없음). 이후 규칙(ul 114)으로 재개에는 사용자 명시 승인이 필요하다. `Astra 단독 7/7' = 파일럿 4/4 + 전체 팔 앞 3/3.}
  \item \textbf{참고:} Astra 단독 실제 편 7개 영상을 파드 /data에 저장(4d54506, 19:21) --- 영상 규칙의 첫 적용(F8).
\end{itemize}

\paragraph{A10. 비용 사건과 유료 승인 규칙}
{\small 09-26 16:45 -- 09-27 00:51 KST}
\begin{itemize}
  \item \textbf{결과:} 사건 목록: (1) E-ACC P1 크레딧 소진 --- 빈 답 37개, 장부 예약 5,043원은 미과금, 과금 1,352.7원(P108); (2) AS 전체 팔 크레딧 HOLD; (3) AS 전체 팔 3편 뒤 비용 지적으로 멈춤(A9); (4) E-OpenVLM 대리의 gpt-5-mini가 7호출(23.3원) 뒤 사용자 지시로 멈춤(결과 \texttt{open\_\allowbreak{}vlm\_\allowbreak{}solo.\allowbreak{}md} 1절, B3).
  \item \textbf{결과:} 지출(책 05): 확정 누적 19,000.4원 = 탐침 6,933.4 + 프롬프트 검진 2,093.9 + Astra 단독 파일럿 3,080.9 + E-ACC 6,892.2; 여기에 E-OpenVLM 대리 1,322.4원을 따로 더한다(책 05 표기). AS 전체 팔 앞 3편 1,359.0원은 그 장부(등록 변경 3).
  \item \textbf{사용자} (ul 114, 09-26 21:55 KST): \intent{돈드는건 내 허락맡고 하기}
  \item \textbf{그래서(결정):} user-log 114(2882208, 21:56): 유료 호출(Astra·OpenAI 등)은 목적·호출 수·예상 비용을 먼저 제안하고 사용자 명시 허락 뒤에만 --- ul 87의 `승인 불필요'를 대체, 자체 검사 의무는 그대로. 가드·등록 문구(99b3b58, 22:23, Task 26), 모든 실제 Astra 클라이언트는 승인 근거 없이 거부·\texttt{run\_\allowbreak{}r5} 기본 mock(fddefd4, 23:59, §85 후속 2), CLAUDE.md 줄(02de6cb, 09-27 00:51). 앞서 \texttt{insufficient\_\allowbreak{}quota} 치명 정지·미응답 요청 별도 보고(688992c, 19:14, Task 21).
  \item \textbf{다음:} 유료 E-Couple(상한 25,000원)·E-Astra-necessity·AS 1단계·자동 맥락 A2 유료 절제·E-AT는 모두 사용자 승인 전 실행 금지.
\end{itemize}

% Generated by D:/tools/scratch_qdd/mmrec/render.py from data_b.py - edit the data, not this file.
\subsection{B. Astra 대체 --- 학습할 수 있는 크기의 상위 모델로}\label{rec:B}
\noindent\textbf{한 줄:} 목표는 학습 가능한 크기의 오픈 VLM이 Astra 자리를 맡는 것(§96). 영샷은 gpt-5.2 대리 0/4·8B 0/5로 인식에서 막혔고, 무료 시뮬 참값으로 LoRA 학습한 8B는 4/4(접근 오차 125.5 $\rightarrow$ 2.6 mm). 최종 상위는 35B, 8B는 임시. 손으로 준 정보 없이도 되는지(E-STRIP8)와 거리를 모델이 아는 두 방법(E-PT)은 등록 뒤 진행 중.

\paragraph{B1. 배경: Jev를 못 쓰게 됨 $\rightarrow$ 로컬 VLM 선택기}
{\small 09-24 18:44 -- 19:10 KST}
\begin{itemize}
  \item \textbf{사용자} (ul 46, 09-24 18:44 KST): \intent{jev는 지금 사용이 불가능한상황임을 확인했어 애초에 텍스트 기반의 api 활성화 기반자체가 우리랑 엮기에는 약간은 부족할 수 있겠단 생각을 해서 jev를 대체할 무언가 방법을 찾아서 해볼 필요가 있을 것 같아}
  \item \textbf{사용자} (ul 47, 19:10 KST): \intent{아스트라써야해 알지?}
  \item \textbf{그래서(결정):} 정본 §44(09-24 18:55 KST): 하위 결정기 = 로컬 VLM typed 선택기 Jev-L(Qwen3-VL, vLLM, 보기 토큰 로그 확률). 사전 시험: 4B 텍스트 + 머리캠 동시 4개 p95 0.261 s. Astra는 위층 계획기로 유지(ul 47 선택).
  \item \textbf{참고:} 이 하위 자리는 이후 융합 VLA(Qwen3-VL-4B + 행동 전문가, §51·§58)로 흡수되었고, 이 줄기(B)는 위층 Astra 자리를 무엇이 대신할지를 다룬다.
\end{itemize}

\paragraph{B2. 선생-학생 $\rightarrow$ 목표 재정의: 학습 가능한 크기로 Astra 대체}
{\small 09-26 19:4x -- 20:00 KST}
\begin{itemize}
  \item \textbf{사용자} (ul 108, 19:4x KST): \intent{선생으로 활약을 잘한다면 꽤 괜찮을거라고 보거든? Llm의 단점을 해결해야해 느린 답변 Vla의 낮은 정확성}
  \item \textbf{사용자} (ul 108, 19:4x KST): \intent{둘중 옳은 방식으로 해야지}
  \item \textbf{그래서(결정):} 정본 §96(047ebc8, 19:57): 시뮬 선생 데이터(E-TEACH-0: 학생 VLA가 시뮬 폐루프를 몰고 선생이 라벨)는 본학습 자리만, 실물 동기 선생 조종(E-TEACH-R)은 추가학습 자리(실물 접근이 없어 계획만). 선행 조건 = 런타임 멈춤 고리 수정(C7).
  \item \textbf{사용자} (ul 109, 20:0x KST): \intent{나는 학습할 수 있는 llm크기에서 아스트라를 대체하는게 목적이야}
  \item \textbf{그래서(결정):} §96 개정 1(ce9c567, 20:00): 주 학생 = 학습 가능 크기 오픈 VLM(Qwen3.5-27B·35B-A3B 급, H200 1--2장 LoRA; 확장 122B-A10B QLoRA), Astra는 증류 선생(E-TEACH-L). 성공 = DEV 짝 40에서 Astra 단독 비열등(여유 3/40), 지연 p50 $\leq$ Astra의 1/3, API 0원. 양은 무료 참값 모델(\texttt{astra\_\allowbreak{}solo/\allowbreak{}truth.\allowbreak{}py}), Astra는 핵심 상태만(단계 $\leq$ 8,000원); 팔 L-T(참값)·L-A(Astra)·L-TA. VLA DAgger는 보조 트랙 E-TEACH-V.
\end{itemize}

\paragraph{B3. 영샷 오픈급 모델로 바로 대체? --- gpt-5.2 대리 0/4}
{\small 09-26 19:45 -- 20:17 KST}
\begin{itemize}
  \item \textbf{해 봄:} 조사 4b91f98(19:45): 목표 오픈 VLM Qwen3.5-397B-A17B(Apache-2.0, 4 H200 FP8로 호스팅 가능), OpenAI 대리 gpt-5.2(같은 Qwen 카드 표 기준). 내려받던 부분 파일 96 GB는 사용자 정정으로 멈추고 삭제. 사전 등록 5c021c6(19:49, 사용자 `허용': gpt-5.2 $\times$ 4편, 상한 3,000원): \texttt{astra-\allowbreak{}solo@v2} 그대로 모델 id만 바꿈, effort low, 조기 종료 20호출/60 s.
  \item \textbf{참고:} 등록 변경 1(f37051c, 20:06): 사용자 정정 --- 목표는 우리가 미세조정할 수 있는 크기(Qwen3.5-27B), 대리 gpt-5-mini 추가. gpt-5-mini는 7호출(23.3원) 뒤 사용자 지시로 멈춤(부분 편, 판정에 안 씀).
  \item \textbf{결과:} 4fbf4e7(20:17): gpt-5.2 0/4, 파지 + 들기 0/4, 무효 0/79. 접근 목표 xy 편 중앙 104--207 mm(같은 판 Astra 2.6--48.9 mm) --- 격자 깊이 방향 x를 12--19 cm 멀게 읽음. 16.4원/호출, 지연 p50 7.2 s. 1,299.1원(+ mini 23.3원).
  \item \textbf{그래서(결정):} 판정 PERCEPTION: 영샷으로는 격자 위치 읽기를 못 한다 $\rightarrow$ 학습 가능 크기 오픈 VLM도 영샷 대체는 기대하기 어렵다; E-TEACH-L의 출발점은 `인식 0', 교사 데이터는 접근 목표 좌표(인식)를 가장 많이 가르쳐야 한다. 같은 인터페이스의 무료 Qwen3-VL-8B 영샷도 0/5(A8 사전 실행).
\end{itemize}

\paragraph{B4. 로봇 자기 정보 말고는 전부 자동화 + 모델 수 최소화 + 깊이 가정}
{\small 09-26 20:3x -- 20:47 KST}
\begin{itemize}
  \item \textbf{사용자} (ul 110, 20:3x KST): \intent{음 이 뭐랄까 로봇의 형상만 보고도 로봇의 엔드포인트를 알아서 알아내고 그것이 행하는 방향을 알아내는걸 자동화 하는걸 할필요가 있어보임 지금 우리가 준 정보라는 것 자체들 전부다 그리퍼 형상, 카메라,범위 같은 로봇이 자기가 스스로 할 수 있는것 외에는 다 자동화를 할 수 있어야해 근데 이게 뭐 여러모델을 이어붙이는건 좋은데 그 모델을 최소화 해야하는 게 좋은 목표인거고}
  \item \textbf{사용자} (ul 110, 20:3x KST): \intent{댑스를 쓴다는걸 가정을 해야겠다}
  \item \textbf{그래서(결정):} 정본 §97(36a7309, 20:47): \texttt{astra-\allowbreak{}solo@v2} 인벤토리 R 16 · P 5 · L 3 · X 3 · K 1. 로봇 자기 정보(R)는 URDF·관절 상태·camera\_info·자기 시험에서 코드가 자동 생성해 계속 준다; 장면(P)은 깊이 + VLM 점 + 코드(RANSAC 평면·평면 위 군집·역투영); 요령(L)은 학습; X는 절제로 판단. 시각 자기 인식은 상위 VLM의 점 출력(시뮬 FK 무료 라벨) + FK 대조. 깊이: 시뮬 렌더 깊이, 실물 머리 ZED Mini 스테레오 계산·손목 D405. 모델 목표 2개(상위 VLM + VLA).
  \item \textbf{다음:} 절제 사다리 A0(무료) $\rightarrow$ A1(무료) $\rightarrow$ A2(유료 $\leq$ 8,000원) --- 모두 [예정], 결과 없음.
\end{itemize}

\paragraph{B5. 가설 H-L --- 상관관계만 학습한 상위 LLM이 더 빠르고 비슷하거나 낫다}
{\small 09-26 21:44 -- 21:50 KST}
\begin{itemize}
  \item \textbf{사용자} (ul 112, 21:44 KST): \intent{사실 실시간으로 아스트라를 키기에는 문제가 좀 있다는 생각이 있어 아스트라가 주역인데 아스트라가 실시간적으로 참여하기에는 너무 길다 감독을 하기에는 시간이 너무 길다는 문제가 있음 따라서 상위 llm이 아스트라보다는 성능이 낮지만 이 엔드포인트와 로봇의 상관관계 이동의 상관관계만을 학습한 것이 더 빠르고 성ㅇ능은 비슷 오히려 더 좋을 것이다가 가정이야}
  \item \textbf{사용자} (ul 113, 21:50 KST): \intent{이런 말두ㅡㄹ도 아직 확정은 아니지만 논문에 미리 업뎃해두고 기억해줘야헤}
  \item \textbf{그래서(결정):} 정본 §96 보충 1(e4156f8, 21:48): H-L 등록, 측정 기준 (a) 상위 호출 지연 p50 $\leq$ 2 s·Astra의 1/3 이하, (b) DEV 짝 폐루프 성공이 Astra 단독과 3/40 이내, (c) 접근 xy 오차 $\leq$ Astra(중앙 약 14 mm)·7--10 cm 오표적 감지 > Astra 0/13. 전제 표 P13. 논문 수시 갱신(054f250, 21:48, 주황 표시).
  \item \textbf{참고:} 당시 근거: gpt-5.2 대리 0/4(104--207 mm), Qwen3-VL-8B 0/5, Astra 자신도 7--10 cm 드리프트 0/13; Qwen3.5-35B-A3B 지연 약 1--2 s는 추정(실측 아님).
\end{itemize}

\paragraph{B6. 35B 전에 8B로 경향부터 --- E-TEACH-L8, 4/4}
{\small 09-26 21:0x -- 23:53 KST}
\begin{itemize}
  \item \textbf{사용자} (ul 117, 21:0x KST): \intent{학습을하면 35b에서 될것이다하는 가설이 진짜 너무 좋긴한데 이게 학습시간이 너무 오래걸리기 때문에 바로 35로가기는 너무 문제가 큼 그래서 8B를 학습을 해서 경향성이 보이는지만 파악을 하면 더 좋을 것 같기는함}
  \item \textbf{해 봄:} 사전 등록 72efee3(21:20), 변경 1 e7960de(21:49, 학습 전: 미세 배치 8 $\times$ 누적 2). Qwen3-VL-8B LoRA SFT, 시뮬 참값 라벨을 \texttt{astra-\allowbreak{}solo@v2} 형식 그대로(DART식 교란 포함). TRAIN 300편 $\rightarrow$ 조종 2,496행(반복 가중 뒤 4,025) + 인식 보조 QA 2,496 = 에폭당 6,521; LoRA r16 2 에폭.
  \item \textbf{결과:} c253ba1(23:48): 오프라인 DEV 305 스냅숏 접근 목표 xy 중앙 영샷 125.5 $\rightarrow$ 미세조정 2.6 mm(대상 위 단계 158.8 $\rightarrow$ 7.2, 첫 호출 163.4 $\rightarrow$ 8.1), 그리퍼 행동 정확도 0.32 $\rightarrow$ 0.95, 유효 JSON 1.00. 폐루프 DEV 4판: 미세조정 4/4(편당 6호출, 성공까지 11.5--11.9 s, 접근 편 중앙 7.7--19.0 mm) 대 영샷 0/4. 학습 1.39 GPU-h(약 4.4k 토큰/s), 지연 p50 3.5 s, 유료 0원.
  \item \textbf{그래서(결정):} 판정 TREND·STRONG(오프라인) $\rightarrow$ WORKS(폐루프). DAgger는 안 함(GPU 조건 초과·이미 천장). 뜻: 영샷 실패는 크기보다 과제 데이터가 없던 문제 쪽; 8B가 이 과제에서 천장이라 35B의 크기 효과는 가르지 못한다. 논문 수시 갱신(35c9fbb, 23:53).
  \item \textbf{다음:} 더 어려운 평가(OOD 높이·물체·방해물, 새 과제, 실물 프레임)와 다양한 데이터(E5·E6), 35B-A3B 처리량 스모크. 한계: n = 4, 참값 라벨(Astra 라벨 아님), 과제 하나, 프롬프트에 손으로 쓴 정보가 있음($\rightarrow$ B9).
\end{itemize}

\paragraph{B7. m 좌표 없이 조종 $\rightarrow$ 점 찍고 행동 + 거리를 모델이 아는 두 방법(E-PT)}
{\small 09-26 23:4x -- 09-27 00:50 KST}
\begin{itemize}
  \item \textbf{사용자} (ul 118, 23:4x--23:5x KST): \intent{로봇 좌표로 x, y, z가 몇 m인지 없이 할방법을 관련 서적 다뒤져서 찾아봐보 ㅏ뎁스여도되고 이미지여도 되고 다괜찮아}
  \item \textbf{사용자} (ul 118, 23:4x--23:5x KST): \intent{로봇 좌표 x, y, z를 m 단위로 말하지 않고 조종하는 방법도 좋지만 어떻게 그걸 추정할 수 있는지에 대한 방법론도 괜ㅊ낳아}
  \item \textbf{결과:} 조사 1d81290(00:01): 상위 VLM은 m 숫자 대신 영상 위에서 고른다(점 $\rightarrow$ 물체 ID $\rightarrow$ 상대 방향), m 값은 코드가 깊이·카메라 자기 정보로 계산. 근거: SPF 같은 백본에서 점 100 \% 대 글 수치 7 \% 대 후보 선택 40 \%; 직접 metric은 SpatialRGPT $\pm$25 \% 성공 41 \%, CA-MLLM 카메라 크기 $\times$ 0.8에 F1 46.5 $\rightarrow$ 25.8. 추천 1 점 + 역투영 + 평면 위 군집, 2 물체 ID 선택, 3 상대 방향 $\times$ 크기 등급 + 손목 서보.
  \item \textbf{사용자} (ul 120, 23:4x KST): \intent{"그리퍼가 어디로 가야 하나"를 이미지 위에서 배우는 방식으로 해봐야지 지금 우리가 못하는게 뭔데 부족한게 뭐야?}
  \item \textbf{사실 확인:} 부족했던 것: (i) \texttt{astra\_\allowbreak{}solo}에 픽셀 목표 명령 없음, (ii) \texttt{astra\_\allowbreak{}solo} 월드가 깊이를 끔, (iii) 우리 데이터에 픽셀 라벨 없음(자체 시뮬은 투영으로 무료), (iv) 공개 MolmoAct식 변환은 다른 에이전트 몫.
  \item \textbf{사용자} (ul 121, 09-27 00:0x KST): \intent{근데 거리정보나 이런거 없이도 되면 너무 더 좋은 결과가 될 수 있을것 같은데 두가지 방법론이야 없게 해서 기존데이터를 그냥 사용하돼 그 거리 정보를 잘 알 수 있게 하던지 아니면 거리를 크리를 알 수 있게 하던지}
  \item \textbf{사용자} (ul 122, 00:1x KST): \intent{2 방법 모두 아는 방법으로 가야겠어 하자 해보자}
  \item \textbf{그래서(결정):} 정본 §97 보충 2 + 사전 등록 18533d6(00:20): 주 팔 ND-1 \texttt{nd-\allowbreak{}xyz}(RGB만, 격자·탁자 높이 없이 xyz --- 암묵)·ND-2 \texttt{nd-\allowbreak{}est}(탁자 높이·물체 위치·크기를 먼저 추정한 뒤 xyz), 비교 팔 \texttt{pt}(점 0--1000 + 높이 의도, 코드가 머리 깊이·보정·평면 위 군집으로 목표)·L8-xyz(격자), 선택 \texttt{nd-\allowbreak{}pt}. F0 영샷 점, F2 같은 L8 편·816걸음 LoRA $\times$ 팔, DEV + OOD-H(탁자 0.82·0.88). 판정 규칙은 결과 전 고정.
  \item \textbf{결과:} 변경 1(fbf45f7, 00:36, 본 실행 전 스모크 DEV 8편): 깊이 해석기에 로봇 자기 가림, 큰 놓을 곳 라벨 허용 25 mm, 관문 G1 = 글 + 라벨 동일(pt 라벨 결측 4/60, v2 요청 60/60·라벨 60/60이 L8과 동일, 머리 PNG 평균 차 0.46/255). 855ffa5(00:50): 평가 세트에 팔 라벨이 없는 상태도 유지(label\_missing).
  \item \textbf{결과:} 변경 2(64bb9d6, 02:06, 미세조정 팔 답은 아직 없음 --- 학습 중): 이미 본 기준 L8-xyz의 OOD-H($\pm$3 cm)는 오프라인 접근 xy 중앙 2.5 mm·접근 3D 중앙 31.2 mm(z가 탁자 변화만큼 어긋남), 폐루프 4/4(패드 4.5 cm가 $\pm$3 cm z 어긋남을 흡수) --- $\pm$3 cm로는 팔을 가르기 어려워 탁자 0.78·0.92 m($\pm$7 cm) 넓은 OOD-H를 둘째 판정으로 추가(주 판정은 $\pm$3 cm 그대로).
  \item \textbf{진행 중·대기:} \tentative{결과 문서(\texttt{results/\allowbreak{}pt.\allowbreak{}md})는 아직 없다 --- 진행 중. 유료 비교(Astra·gpt-5.2 $\times$ pt/nd-est)는 제안만, 사용자 허락 전 실행 금지.}
\end{itemize}

\paragraph{B8. 최종 상위 = 35B, 8B는 임시 대리}
{\small 09-27 01:2x -- 01:30 KST}
\begin{itemize}
  \item \textbf{사용자} (ul 125, 01:2x KST): \intent{상위를 8B로 하는건 임시인거고 35B짜리는 학습하는게 너무 커야하니까. 최종은 35B가 목적인거야. ㅇㅋ?}
  \item \textbf{그래서(결정):} 정본 §96 보충 2(ce7ba33, 01:30): 최종 상위 = Qwen3.5-35B-A3B, 8B(E-TEACH-L8)는 35B를 바로 학습하기엔 너무 커서 경향만 먼저 보는 임시 대리. 같은 보충의 의도 공유 프로토콜은 C9.
\end{itemize}

\paragraph{B9. 손으로 쓴 프롬프트 정보가 추론 때도 필요한가 $\rightarrow$ E-STRIP8}
{\small 09-27 01:2x -- 01:35 KST}
\begin{itemize}
  \item \textbf{사용자} (ul 126, 01:2x KST): \intent{저 프롬프트는 누가 만들어준건데? 저거는 학습할때는 들어갈만한데 추론시에도 있었다는거?}
  \item \textbf{사실 확인:} \texttt{astra-\allowbreak{}solo@v2} 프롬프트는 우리(Claude 에이전트)가 손으로 썼고, E-TEACH-L8은 학습·추론 모두 이 프롬프트를 썼다: 탁자 높이 0.850, 물체 모양·크기, 그리퍼 기하, 잡기·놓기 요령, 카메라 자세, 작업 상자, 탁자 높이에 그린 격자·낙하선.
  \item \textbf{해 봄:} 사전 등록 5511d5d(01:32, 학습 전, 무료): 팔 S-full(L8 그대로)·S-min(학습·추론 모두 최소 요청: 탁자 높이·격자·낙하선·물체 크기·요령 뺌, 로봇 자기 정보는 남김)·S-drop(학습 때 특권 정보 드롭아웃, 추론은 최소)·S-full$\rightarrow$min(L8 어댑터에 최소 요청)·S-stale(OOD-H에 손으로 쓴 0.850). 같은 L8 상태·816걸음, DEV 305 + OOD-H 오프라인, 판정 NOT\_NEEDED / NEEDED / UNCLEAR은 결과 전 고정. \texttt{strip(v2, \{table, overlay\})} = E-PT \texttt{nd-\allowbreak{}xyz@v1}와 바이트 동일(시험).
  \item \textbf{결과:} 변경 1(8700871, 01:35, 학습·평가 전): G-stale 기준 = 글 동일 + 영상당 다른 화소 $\leq$ 100(재생성 142장 중 48장이 서브 mm TCP 차로 낙하선 1--69 화소가 다름).
  \item \textbf{진행 중·대기:} \tentative{학습·평가 전 --- 결과 없음.}
\end{itemize}

\paragraph{B10. 35B-A3B 바로 할 수 있게 준비만 --- 지연 약 1 s (준비 측정, 과제 성능 아님)}
{\small 09-27 01:0x -- 02:4x KST}
\begin{itemize}
  \item \textbf{사용자} (ul 135, 01:0x KST): \intent{35B-A3B는 일단 바로 할 수 있게 준비만 다 해둬보자}
  \item \textbf{해 봄:} 준비만(본 학습 안 함, 1e4ed20): \texttt{Qwen/\allowbreak{}Qwen3.\allowbreak{}5-\allowbreak{}35B-\allowbreak{}A3B}·FP8판 내려받기, 서빙 지연 실측, 학습 스모크 50걸음, 병합-서빙 왕복, 실행 스크립트 \texttt{tools/\allowbreak{}teach\_\allowbreak{}35b/\allowbreak{}train.\allowbreak{}sh}, 사전 등록 초안 \texttt{docs/\allowbreak{}stage3/\allowbreak{}prereg\_\allowbreak{}teach\_\allowbreak{}35b.\allowbreak{}md}(한계 찾기 지표 3.1 포함).
  \item \textbf{결과:} \texttt{results/\allowbreak{}teach\_\allowbreak{}35b\_\allowbreak{}ready.\allowbreak{}md}: FP8·H200 한 장·한 요청씩·생각 끔 영샷 --- solo p50 1.11 s·p95 1.26 s(유효 59/60), couple@v2 p50 1.50 s; 생각 켬은 6,000토큰까지 잘려 쓸 수 없음. 학습 스모크(fla 커널) 6.07 s/걸음 $\rightarrow$ L8 2 에폭 약 1.4 GPU-h(8B와 같음). 병합 $\rightarrow$ vLLM 왕복 DEV 20요청 모두 유효. 유료 0원.
  \item \textbf{그래서(결정):} 학습·평가 기본 경로는 solo 형식(ul 134·§98, C11) --- 결합 형식 지원은 코드에 남긴다.
\end{itemize}

\paragraph{B11. 명령어 범위 --- 집고 놓기 먼저, 한계 확인 뒤 확장}
{\small 09-27 02:5x KST}
\begin{itemize}
  \item \textbf{사용자} (ul 136, 02:5x KST): \intent{오픈데이터의 경우엔 엄청 다양한 테스크도 있는거지?}
  \item \textbf{사용자} (ul 136, 02:5x KST): \intent{이게 맞지}
  \item \textbf{그래서(결정):} 상위 명령어는 당분간 지금 형식(손끝 목표·edit·그리퍼·stop, 위에서 잡기)으로 두고, 공개 데이터의 다양한 과제(RobotisSW 322종, BEHAVIOR 50과제, RoboTwin 50종, Robo2VLM 3,396과제 등)는 인식·판단 학습에만 쓴다. 상위 단독의 한계가 확인되면(§98) 손목 회전·접근 방향·밀기·당기기로 명령어를 넓힌다. 논문 3.2절에 한 문장으로 반영.
\end{itemize}

% Generated by D:/tools/scratch_qdd/mmrec/render.py from data_c.py - edit the data, not this file.
\subsection{C. 의도 공유 프로토콜과 결합 --- 상위는 구간과 할 일을, VLA는 그 안의 `언제'를}\label{rec:C}
\noindent\textbf{한 줄:} VLA를 흔한 VLA처럼 쓰지 않고 말단 보정기·세부 조종 스틱으로 두는 틀(§87·§90·§91·§92)을 세우고 결합 배선(Task 1--26)을 끝냈다. 무료 사전 실행에서 VLA 단독이 0/12였고 그 중 8편은 런타임 교착이었다(수정함). 상위$\leftrightarrow$VLA 의도 공유는 필수(§96 보충 2). 남은 것은 학습 --- ser-A-min-3 VLA 재학습, 결합 형식 상위 학습, 결합 폐루프 검증.

\paragraph{C1. 틀: VLA를 VLA처럼 쓰지 않는다}
{\small 09-25 23:18 -- 09-26 01:07 KST}
\begin{itemize}
  \item \textbf{사용자} (ul 75, 09-25 23:18 KST): \intent{그리고 한가지 궁금한게 이게 VLA를 따르기는 하는데 VLA는 일반적으로 VLM모델을 엔드투 엔드학ㄱ습을할때 이 엔드포인트를 어디로 움직여야한다는 모델로서의 기능이 아니라 테스크 전체에 활용하는 vlm이잖아? 근데 우리는 이 vlm의 역할을 단순히 오락실 게임기처럼 위아래 앞뒤와 같은 것중에 선택을 하게 하는것에 활용을 한다는 점임을 강ㅇ조를 해야할것 같고 실제로도 그렇게 활용될 수 있도록 하고 있는거지?}
  \item \textbf{사용자} (ul 77, 23:50 KST): \intent{내가 기존의 VLA 모델에서 문제라고 봤던점은 지금 내가 하는 행동 의 단위 자체가 모든 동작들마다 일관적일 수 없다는 점이야. 생각을 해보자. 우리가 옷을접을때 그냥 책상 위에 있을때는 잘 접을 거야. 그런데, 책상위에 벽돌 30개가 있어서 옷을 잡으러 갈때 계속 방해해버려 VLA는 이걸 절대 못풀어. 그이유는 VLA는 행동의 연속성을 의미단위로 풀려하지만 그 행동들은 의미를 가지기에는 너무 짧은 구간들만 보고 있는거거든 볼거면 전체의 흐름을 모두 본다음에 그 의미와 행동을 매칭해야하는데 그게 아니잖아. 즉 이행동을1}\textasciitilde\intent{3초간의 이미지만 가지고 내가 지금 뭐하고 있는거게}\textasciitilde\intent{ 하는 거나 다름이 없는 이상한 거야. 그래서 효율적인 학습을 위해선 이 전체적인 학습의 흐름이 뭐다 라는걸 아스트라가 잡아주는 과정이 굉장히 중요하다고 생각하는거야. 이걸 어떻게 유기적으로 활용할 수 있을까? 나랑 지금 브레인 스토밍해보자}
  \item \textbf{사용자} (ul 77, 23:50 KST): \intent{내가 방해물을 예시로 들었지만 그게 아니라 그냥 다른환경에 취약하다는 얘기야}
  \item \textbf{결과:} 가설 보정 조사(a8d4242, 09-26 00:10): 좁은 시연으로 순진하게 끝단 학습하면 국소 지름길을 배워 새 위치·시점에서 무너지는 경향은 체제 한정 지지; `무조건 과적합'은 반박; `짧은 창이 원인'은 근거 없음. 조향 표현 조사(1079716): 동작 수준 명령 인터페이스는 지지, `어디로' 공간 목표는 없음.
  \item \textbf{사용자} (ul 81--82, 09-26 00:46--00:50 KST): \intent{아게 진짜로 잘 모르긴하겠다 아스트라를 어케 쓰지?}
  \item \textbf{사용자} (ul 81--82, 09-26 00:46--00:50 KST): \intent{나는 VLA를 VLA 처럼 쓰고 싶지 않은거야 뭔소리인지 알아? 단점이 있기에 그건 순간 순간 어디로 가야하는지 에ㅐㄴㄴ드포인트 정밀 보정하는 역할을 하는 용도로 쓸거여}
  \item \textbf{사용자} (ul 81--82, 09-26 00:46--00:50 KST): \intent{그래서 전체적으로 다 엔드포인트를 위한 움직임 만을 하난 파이프 라인인거야. 이해했어? 근데 그럼 계획을 VLA로 부터 뺐어 왔으니까 이걸 아스트라가 담당해야하는데 아스트라가 잘하는 걸 하게 해야한다는거지}
  \item \textbf{그래서(결정):} 논문 중심 틀(cd5210f, 01:07 논문 갱신): VLA = 순간순간의 말단 정밀 보정기, VLM은 오락실 조이스틱처럼 typed 보기만 고름; 과제 이해·흐름·계획은 위에서(Astra J1 + 흐름 교정, A1·A2).
\end{itemize}

\paragraph{C2. 결합 설계 승인, MolmoAct를 주 참고로, 구현 계획}
{\small 09-26 01:44 -- 03:00 KST}
\begin{itemize}
  \item \textbf{사용자} (ul 84, 01:44 KST): \intent{몰모 엑트를 많이 참조분해해서 잘사용 적용해보면 좋을듯}
  \item \textbf{사용자} (ul 85, 02:11--02:13 KST): \intent{1}\textasciitilde\intent{3 다 허용할게 몰모엑트 많이 참고하고}
  \item \textbf{사용자} (ul 85, 02:11--02:13 KST): \intent{1,2,3 허용했지만 옳은 방향으로 바꿔도돼 확인해보고}
  \item \textbf{그래서(결정):} 정본 §84(c915535, 02:11): 결합 설계 승인, E-MA1 승인, 같은 진단 반복 시 falsify 먼저 평가, MolmoAct 채택 순위(미래 손끝 궤적 보조 $\rightarrow$ 이력 덧그림·궤적 조건 $\rightarrow$ 층별 KV $\rightarrow$ 깊이 보조 보류). 보충 1(a5816e1, 02:13): falsify 재사용은 같은 에피소드·같은 실패 서명에서 한 번까지, E-MA1 해석 주의(일반화 근거로 쓰지 않음), 호출 빈도는 비용 한도로.
  \item \textbf{그래서(결정):} 구현 계획(bb1f5f7, 02:59)과 메인 판정(7bb9035, 03:00): 두 층 합의 + T1 전제, 탐침 예산 15,000원, 실행은 R7 E2E 준비 기준점 뒤 subagent-driven.
  \item \textbf{다음:} R7 E2E 준비 기준점 도달(36·37회차 연속 무결, 태그 \texttt{stage3-\allowbreak{}e2e-\allowbreak{}ready} = 3d4738a, 8de6b80 09:46) 뒤 구현 착수(C5).
\end{itemize}

\paragraph{C3. 잡기·놓기 판단은 VLA의 선택지, VLA 자체가 세부 조종 스틱}
{\small 09-26 11:14 -- 11:17 KST (구현 15:17)}
\begin{itemize}
  \item \textbf{사용자} (ul 93, 11:14 KST): \intent{근데 그냥 소능 움직일때 물체가 근접하지 않았다면 상관 없지만 물체에 가깝게 다가가서 미세조정이 필요한데 아스트라의 반응성이 너무 느리면 그걸 보조하는 역할이라서 어떤걸 잡아야한다 이때 놓아야한다등에 대한 지식을 가져야하는데}
  \item \textbf{사용자} (ul 94, 11:17 KST): \intent{그 잡고, 놓고 하는 판단 자체도 vla가 선체택지에 있어야한다는 뜻}
  \item \textbf{사용자} (ul 94, 11:17 KST): \intent{vla자테가세부 조종스틱이기도하지 물론}
  \item \textbf{그래서(결정):} 거리에 따른 권한 분배(a4ae36e): 먼 이동 구간은 조이스틱·Astra를 강하게 따르고, 접촉 근처(\texttt{core.\allowbreak{}near\_\allowbreak{}contact} $\leq$ 5 cm)는 VLA 자기 과제 지식이 주도. 정본 §87(6e9f778): 융합 VLA typed 결정에 그리퍼 질문(쥐기 close / 놓기 open / 유지 keep); 선택지는 `의도', 실행 허가는 두 층 관문 + T1. 계층(6127f92): Astra = 느린 큰 조종, VLA 결정 = 세부 조종 스틱(0.33 s), 행동 전문가 = 그 입력을 연속 움직임으로.
  \item \textbf{결과:} 구현 확인(Task 12, 340c5dc 15:17): 그리퍼 질문은 어디에도 없었고 \texttt{phase}는 쥐기·놓기 의도를 명시하지 않았다 $\rightarrow$ 그리퍼 질문을 융합 VLA에만 추가, 라벨 규칙은 \texttt{harvest/\allowbreak{}intent.\allowbreak{}py} 한 곳(R2 = 기록 단계, S-E2E = 그리퍼 명령 속도 0.176 /s).
\end{itemize}

\paragraph{C4. 구간 계획(§90) · 도착 시 대조(§91) · 구조 확정(§92)}
{\small 09-26 14:37 -- 14:48 KST}
\begin{itemize}
  \item \textbf{사용자} (ul 100, 14:37 KST): \intent{상위 아스트라가 어디로 어케 움직일거고 나는 지금 다가가는 행동ㅈ이다, 이 구간선 잡는걸 할거다. 이구간선 드는걸할거다 하는ㅇ계회도 담고 vla는 그 구간이 길기에 그구간에서 잡는게 있다곤 하는데 언제 잡을지를 상황을보고 결정하는 겆ㄱ}
  \item \textbf{그래서(결정):} 정본 §90(aed65ff, 14:37): Astra 답에 구간 계획(지금 구간·이 구간에서 할 일·다음 구간), VLA는 구간 의도를 입력으로 받아 그 안에서 `언제'를 스스로 정함(공간 목표는 여전히 없음).
  \item \textbf{사용자} (ul 101, 14:40 KST): \intent{그리고 아스트라의 판단이 몇초전 이미지와 상황을 보고한거지 그리고 실제vla가 조정한건 어땠는지의 비교를가지고도 추가판단이 가능해야함}
  \item \textbf{그래서(결정):} 정본 §91(5158cbc, 14:40): 도착 시 대조 --- 판단 기준 시각의 상태와 지금, Astra 명령과 그사이 VLA의 실제 움직임을 비교해 이미 됨 / 아직 유효 / 상황 바뀜 / 충돌을 판정, 다음 요청 문맥과 추가 학습 기록에 넣음.
  \item \textbf{사용자} (ul 102, 14:44--14:48 KST): \intent{이제 어떻게 서로 엮어서 최대 효율을 내게끔 파이프라인 구조 확정해서 이미지좀 만들어줘봐봐}
  \item \textbf{사용자} (ul 102, 14:44--14:48 KST): \intent{아직 돌고 있는 작업들은 확정이 나야지 파이프라인 확정이지? 아직 확실하지 않은 전제들이 좀 있지?}
  \item \textbf{그래서(결정):} 정본 §92(c410ec2, 14:44): ① Astra(low) 직렬 흐름 ② 도착 시 대조 + 합의 + 권한 a ③ VLA 결정 = 세부 조종 스틱 ④ 행동 전문가 ⑤ 안전 관문 ⑥ 실행 ⑦ 기록 + 그림. 보충 1(553a8ad, 14:48): `배치' 확정일 뿐, 전제 P1--P12(뒤에 P13 H-L)가 검증되기 전까지 잠정.
  \item \textbf{참고:} 뒤 결과로 본 전제(정본 표의 상태 칸은 작성 당시 값): P1 실데이터 준수 0.8 --- E-SR1d 0.536으로 불성립, E-SR1e 진행 중(D7·D8); P3 결합 > VLA 단독 --- VLA 단독 0/12로 짝 비교 보류(§92 보충 2, C7); P4 명령 정확도 --- 수정 방향 13/13, 7--10 cm 알아채기 0/13(A6·A7); P10 E-Astra-necessity --- 초안, 유료 금지; P13 H-L --- 8B 경향 4/4(B6), 35B 미정.
\end{itemize}

\paragraph{C5. 결합 배선 구현 Task 1--26과 리뷰에서 잡은 결함}
{\small 09-26 10:02 -- 23:59 KST}
\begin{itemize}
  \item \textbf{해 봄:} 계획 \texttt{docs/\allowbreak{}superpowers/\allowbreak{}plans/\allowbreak{}2026-\allowbreak{}09-\allowbreak{}26-\allowbreak{}astra-\allowbreak{}vla-\allowbreak{}coupling.\allowbreak{}md}, subagent-driven, 통제자 판정. Task 1 파라미터·답 스키마·\texttt{astra-\allowbreak{}couple@v1}(205db33 10:02), 2 답 게이트(7c56331), 3 가격표·공유 장부 80 \% 정지(0e480ac), 4 직렬 흐름·모의 Astra(e6f163c), 5 두 답 이력(e446070), 6 부드러운 편향(b5c059b), 7 두 층 비가역 관문(T1 근거만)·VLA 빠른 검증(f9eb416), 8 Astra 전용 덧그림·카메라 모델(8350c7e·85bd2ed), 9 CoupleDriver(dd3ce0c), 10 OursRuntime 배선(dd1afbc), 11 falsify 먼저 재사용(8074683), 12 직렬화 \texttt{ser-\allowbreak{}A-\allowbreak{}min-\allowbreak{}3}(340c5dc 외), 13 로컬 VLM 어댑터(58eed4b), 14 폐루프 결합 팔·유료 가드(ea98034), 15·16 E-Couple·E-Astra-necessity 등록 초안(2a99db6·2eea17b), 17 편향 $\times$ 권한 a(b0ee916), 18 \texttt{astra-\allowbreak{}couple@v2}(8c9064a), 19 도착 시 대조(20d735c), 20 프롬프트 검진 가드(a913020), 21 사건 불응기·미응답 보고·quota 치명 정지(688992c), 22 확신도 기록 훅(5a940d2), 25 교착 수정(c835476), 26 유료 가드 ul 114(99b3b58), 최종 검토 수정(fddefd4 23:59).
  \item \textbf{결과:} 리뷰에서 잡아 고친 결함(통제 판정): 구간 줄이 그리퍼 정답을 알려 주던 것(F12, bce0bd8 --- 구간이 사건을 품게, 학습 때만 unknown 구간 p 0.3); 권한 a가 떨어질 때 남은 편향이 1.92 cm·0.36 rad 튀던 것(F17, badb1e2); 실행 청크가 없을 때 결정 중심으로 거짓 화살표를 그리던 것(F18, 9328064); 늦은 답이 반영되던 것·모듈형 스킬이 편향을 매 틱 지워 가짜 충돌을 내던 것(F19·F19b·F19c, ffd7919·3096b26·c326d45); 탈출 청크가 버려지던 것 등(F25, 8bfe694).
  \item \textbf{그래서(결정):} 정본 반영: §84 보충 2 구현 기록(Task 21), §84 보충 8 구현 기록(Task 17: 명령 속도에 a를 곱한 뒤 가속·속도 상한, a = 0이면 Astra 수정 효과 0), §90 구현 기록 1--3, §91 구현 기록(우선순위 판정, 늦은 답은 늘 폐기, 예상 상태 = 지금 손끝 + 실행 청크 변위), §4·§58·§64 보충(C8). §6 결정 투영은 구현하지 않음.
\end{itemize}

\paragraph{C6. 무료 결합 사전 실행(셰이크다운) --- 배관은 되고 VLA가 바닥}
{\small 09-26 15:06 -- 16:28 KST}
\begin{itemize}
  \item \textbf{해 봄:} 사전 등록 9322853(15:06, P3 판정 없음, 0원): VLA = E-SR1c C1, Astra 대역 = Qwen3-VL-8B + \texttt{astra-\allowbreak{}couple@v1}(Astra 비슷한 지연 + 주입 시간 초과), DEV 0--5 $\times$ \{standard, dr\} $\times$ \{VLA 단독, 결합\}, 메인 GPU 1. 변경 1(f9903d7): 분석 지표만 추가.
  \item \textbf{결과:} bd3a8ba(16:28): 24/24편 충돌 0, 동시 1개 위반 0, 시간 초과·늦은 답 5(주입 6), stale 0, 스키마 오류 3/55, 편향 속도 최대 0.0084 $\leq$ 0.08 m/s, 가속 0.32 = 상한, 장부 0원. 성공 VLA 단독 0/12·결합 0/12.
  \item \textbf{결과:} 버그: B1 그리퍼 닫기를 풀 때 한 틱 14 mm 계단; B2 \texttt{max\_\allowbreak{}inflight} 기록이 증거가 못 됨; B3 편향·관문 로그 누락; B4 권한 a·§91·§90·움직임 줄이 이 판에 없음(a\_priv = 0에서 편향 0.42 mm 적용); B5(차단) VLA 단독 0/12 바닥 --- 짝 차가 정의상 0; B6 \texttt{b\_\allowbreak{}contradict} 사건 892회로 선택 단계 쉼이 안 켜짐; B7 편 끝 미응답 요청의 예약 청구; B8 Qwen 스키마 오류 5.5 \%·\texttt{grasped} 과대 주장 15.
  \item \textbf{그래서(결정):} B1--B3는 b0ee916(Task 17), B6·B7은 Task 21, B4는 Task 12·17--19로. B5 $\rightarrow$ 유료 E-Couple 전에 VLA 단독 폐루프 온전성 관문을 둔다.
\end{itemize}

\paragraph{C7. VLA 단독 0/12 분해 $\rightarrow$ E-VLA-solo: 의미 없음, 느리게 해도 안 됨, 교착 확인}
{\small 09-26 16:3x -- 21:22 KST}
\begin{itemize}
  \item \textbf{사용자} (ul 107, 16:3x KST (통제자 경유)): \intent{Vla는 혼자서 의미가 있는지도 보긴해야해 그리고 로본의 동작 속도자체도 조금 느리게 가져가도 좋을듯}
  \item \textbf{결과:} 진단(d752509, 18:44, 기존 기록만·CPU): 12편 중 8편이 런타임 교착 --- T1 모순 $\rightarrow$ M4 hold $\rightarrow$ 결정 비움 $\rightarrow$ 청크 버림 $\rightarrow$ 같은 모순. 방아쇠는 approach 중 청크 그리퍼 폭 77--80 mm가 열림 문턱 80.5 mm 바로 아래(3편), 들기 뒤 쥠 판정이 풀림(5편). 나머지 4편은 VLA 자체 실패(머그를 쳐서 탁자 밖 2, 머그 위에서 안 내려감 2). approach 중 모델 방향 답이 머그 쪽인 비율 36 \%(우연 수준).
  \item \textbf{해 봄:} E-VLA-solo 사전 등록(d752509, 18:44, 무료): E-SR1c C1, 7팔 $\times$ DEV 0--5 $\times$ standard·dr = 84편 --- C5·C3, 재생 1.0·0.75·0.5배, 계획기 규칙·제자리 기준선.
  \item \textbf{결과:} 21a5284(21:22): VLA 다섯 팔 모두 0/12, 계획기 규칙 12/12, 제자리 0/12. 지나침 87.7 $\rightarrow$ 52.9 $\rightarrow$ 17.9 mm(1.0·0.75·0.5배). C5 영구 hold 8/12 대 C3 0/12. C3는 들기 이상 10--11/12까지 가지만 놓기 성공 0. C3 관절 명령 계단(> 0.04 rad/틱) 103--135틱.
  \item \textbf{그래서(결정):} Q1 NOT\_MEANINGFUL(이 체크포인트의 VLA 단독은 성공을 못 냄) · Q2 NOT\_ADOPTED(재생 1.0배 유지, 정본 §92 보충 2 `속도 낮추기는 검증 후 채택') · Q3 CONFIRMED(교착 수정 필수). P3 짝 비교 보류, 유료 E-Couple 금지 유지.
\end{itemize}

\paragraph{C8. 교착 수정과 막힌 파드 검증}
{\small 09-26 21:57 -- 09-27 KST}
\begin{itemize}
  \item \textbf{해 봄:} Task 25(c835476, 21:57): T1 열림 히스테리시스(들어가기 80.5 mm, 나가기 75 mm), hold 3.0 s 뒤 한 스텝 탈출, hold 중 융합 청크 재생, VLA가 쥔 물체의 lift·carry 쥠 해제 $\rightarrow$ object\_lost 실패 경로, 융합 청크 전환 팔 틱당 0.04 rad 클램프. 수정 F25(8bfe694, 22:22): 탈출 청크 재생, 탈출 3연속 $\rightarrow$ \texttt{hold\_\allowbreak{}stuck} 실패, 힘만 떨어지면 \texttt{grip\_\allowbreak{}force\_\allowbreak{}low} 기록만(패드가 50.1 mm 아래로 닫혀야 object\_lost). 최종 검토 fddefd4.
  \item \textbf{진행 중·대기:} \tentative{파드 재측정(Task 25 C)과 결합 사전 실행 r2(Task 24)는 막혀 있다: HEAD의 \texttt{ser-\allowbreak{}A-\allowbreak{}min-\allowbreak{}3}이 기존 체크포인트를 모두 거부(serializer\_ok·last\_step\_ok·files\_ok). [결정 필요] ser-A-min-3 재학습 vs 옛 판 이식 vs 배선만 확인(handoff §0) --- 결합 학습 계획의 추천은 VLA-sim 재학습 먼저(C10 D2).}
\end{itemize}

\paragraph{C9. 상위$\leftrightarrow$VLA 의도 공유 프로토콜은 필수 핵심}
{\small 09-27 01:2x -- 01:30 KST}
\begin{itemize}
  \item \textbf{사용자} (ul 125, 01:2x KST): \intent{저 VLA는 0.33초 마다 스스로 움직여서 공백을 매우지만 상위에서 준 나는 이움직임동안 이러이런걸할거야가 같이 내려오면 그걸 기준으로 실행을 하는거야. 뭔소리인지 이해했지? 행동을 상위가 내 뱉을 때, 나는 이런이런 행동으로 할 생각이야 나는 물건을 향해 다다가는거야 하면서 행동 뱉고 그럼 그사이를 vla가 조종하고 그다음에 이번엔 잡을거야 하면 상위는 거기까지가고 vla는 그걸보고 이번엔 잡는구나 그래 다다가자 근데 언제 잡지?를 vla가하고 상위는 그걸보고 잡았구나 안잡았구나 이런 판단을 더 정확히 할 수 있으니까 그걸 하고. 이해했지? 이건 필수야 꼭 중요한 핵심이야}
  \item \textbf{그래서(결정):} 정본 §96 보충 2(ce7ba33, 01:30): 상위가 내는 모든 행동은 그 구간의 의도(§90 now/do/next)와 함께 내려온다; 답 사이(0.33 s)는 VLA가 그 의도를 기준으로 공백을 메운다; 쥐는 `언제'는 VLA 소유(§87); 상위는 VLA가 실제 한 일을 보고 잡았다/못 잡았다를 더 정확히 검증(§91·T1·V1h).
  \item \textbf{그래서(결정):} 학습 데이터에 대한 결과: 상위(35B)의 학습·추론 출력은 의도/구간 필드와 검증 판정을 반드시 포함 --- \texttt{astra-\allowbreak{}couple@v2} 같은 결합 형식. \texttt{astra-\allowbreak{}solo} 형식 실험(E-TEACH-L8, E-STRIP8, E-PT/ND)은 인식·조종 능력만 잰다.
\end{itemize}

\paragraph{C10. 결합 학습 계획 --- 엮을 부품은 있고, 모자란 것은 학습}
{\small 09-27 01:3x -- 01:39 KST}
\begin{itemize}
  \item \textbf{사용자} (ul 127, 01:3x KST): \intent{근데 이걸하려고 이미 과거에도 좀 정리를 했던게 있었을거야. 얼마나 그걸 잘 엮을 수 있을지 학습은 어떻게 할지가 이제 좀 제대로 잡혔지?}
  \item \textbf{결과:} \texttt{docs/\allowbreak{}research/\allowbreak{}coupled\_\allowbreak{}training\_\allowbreak{}plan\_\allowbreak{}2026-\allowbreak{}09-\allowbreak{}27.\allowbreak{}md}(a54170f, 01:39, 계획만·0원·GPU 0): 런타임 배선은 Task 1--26으로 끝. 없는 것 셋 --- (1) \texttt{ser-\allowbreak{}A-\allowbreak{}min-\allowbreak{}3} 형식으로 학습된 VLA 체크포인트 0개, (2) 결합 형식 상위 학습 데이터·참값 라벨러 \texttt{couple-\allowbreak{}truth} 없음(L8·E-PT·E-STRIP8은 모두 solo 형식), (3) VLA 끝-끝 학습·결합 폐루프 검증 없음. 어긋남 G1(상위 구간 9개 $\rightarrow$ VLA 줄 4개, 의도된 설계)·G2(결합 형식엔 공간 목표 없음)·G3(명시적 검증 필드 없음)·G4(상위 요청에 T1 판정 글 없음)·G5(VLA 학습에 늦은·틀린 의도 없음)·G6(체크포인트 0)·G7(§86 상수는 v1에서 잰 값)·G8(E-AT \texttt{quality:} 줄 없음).
  \item \textbf{그래서(결정):} 정본 §96 보충 3·4: 무료 사다리 V0 VLA 재학습 관문 $\rightarrow$ V1 VLA 단독 $\geq$ 4/12(영구 hold $\leq$ 1/12) $\rightarrow$ V2 결합 사전 실행 r2 $\rightarrow$ V3 E-TEACH-L8C 오프라인(\texttt{upper-\allowbreak{}couple@v1} = v2 + verify) $\rightarrow$ V4 결합 폐루프(학생 8B + VLA 대 8B 단독 동기 대 VLA 단독, DEV 40 짝) $\rightarrow$ V5 35B; V6 유료 E-Couple은 사용자 승인 뒤만. 자리 규칙 §88·§89 그대로. VLA 쪽 E2E는 D12.
  \item \textbf{진행 중·대기:} \tentative{사용자 결정 3개(추천): D1 학생 형식에 \texttt{verify} 필드 추가(예), D2 GPU 순서 --- ser-A-min-3 VLA-sim 재학습을 가장 먼저(예), D3 VLA 구간 줄은 지금 4단계 유지하고 V4 쥐기 시점 오차로 다시 판단. book 01의 VLA-R3·E-VLA-E2E·E-TEACH-L8C·E-COUPLE-8B·E-TEACH-35C는 모두 [예정].}
\end{itemize}

% Generated by D:/tools/scratch_qdd/mmrec/render.py from data_d.py - edit the data, not this file.
\subsection{D. VLA --- 세부 조종 스틱이 결정을 따르게, 그리고 끝-끝으로}\label{rec:D}
\noindent\textbf{한 줄:} 움직임 줄만 채택(+0.032)되고 MolmoAct식 보조·카메라 3대·층별 KV·명령 입력은 모두 불채택. 결정 토큰을 청크가 따르지 않는 문제(E-SR0 0.34--0.43)는 시뮬에서 먼 구간 반사실 분기로 0.991까지 풀었지만 실데이터는 0.536(E-SR1e 진행 중). 확신도는 기록만. 다음은 ser-A-min-3 재학습과 끝-끝 학습(E-VLA-E2E).

\paragraph{D1. 배경: VLM 사용법, 소규모 끝-끝 준비, 하나로 융합}
{\small 09-24 23:05 -- 09-25 23:05 KST}
\begin{itemize}
  \item \textbf{사용자} (ul 54, 09-24 23:05 KST): \intent{어떤식으로 vlm을 써야할지를 좀 생각해봐야할듯 엔드투 엔드 목적으로 사용하게 될거니까 하긋ㅂ할때 어케하게될지}
  \item \textbf{사용자} (ul 54, 09-24 23:05 KST): \intent{A$\rightarrow$B 단계적 (권고)}
  \item \textbf{그래서(결정):} 정본 §51: 먼저 결정층(보기 토큰)만 학습, 뒤에 같은 백본에 행동 전문가를 붙여 전체 끝-끝으로. §55 단계 A SFT: A\_SFT $-$ A\_zero +0.504 [+0.498, +0.511], p95 0.307 s.
  \item \textbf{사용자} (ul 56, 09-25 02:13 KST): \intent{사실 이건 실험인거고 제대로 된 데이터 를 제대로 넣은 실험은 아닌거 아는거지? 지금한건 소규모 테스트로 실제 동작하는지를 본거지?}
  \item \textbf{그래서(결정):} §56: §55는 파이프라인 확인용 소규모 시험 --- 단계 3 수치 대부분에 같은 성격 표시.
  \item \textbf{사용자} (ul 59--61, 02:55--02:59 KST): \intent{그냥 계속 다 자동으로 해주면돼 아직 엔드투엔드는 아니고 모든 파이프라인 전체를 엔드투엔드 학습할 준비까지 모두 준비를 마쳐주면 되는거야 알겠지? 검증해보면서 계속 자동으로}
  \item \textbf{사용자} (ul 59--61, 02:55--02:59 KST): \intent{각 모듈별로 따로 놀면서ㅗ 가각의 성능에 의존하는 파이프라인보다는 최대한 하나로 융합되는게 좋은 연구인거알지?}
  \item \textbf{사용자} (ul 59--61, 02:55--02:59 KST): \intent{어쩔수없는건 물론 그냥 하긴해야하구}
  \item \textbf{사용자} (ul 59--61, 02:55--02:59 KST): \intent{이렇게 모듈 각각 검증해보고 다하고나면 한번 돌아다니면서 객관 검증해보고 2연속 사이클 도는데 문제가 없으면, 이제 다음인 소규모 엔드투 엔드 준비해서 데이터 오픈소스 된것중에 해서 한 10gb정도로만해서 소규모 엔드투엔드 학습한번해보자.}
  \item \textbf{그래서(결정):} §58: 런타임은 하나의 융합 모델(Qwen3-VL-4B 결정 + 흐름 행동 전문가). R7 객관 검증 순회 $\rightarrow$ 2회 연속 무결이 E2E 준비 기준점.
  \item \textbf{사용자} (ul 66, 04:18 KST): \intent{직므은 단순이니까 10으로 하되 나중에는 우리가 데이터하고 할떄는 30으로 모을거}
  \item \textbf{결과:} S-E2E(ROBOTIS 공개 실물 RB1 + RB2, 원본 8.95 GB, 10 Hz $\rightarrow$ 42,813행, train 37,484 / val 1,799): 학습 PASS(bc1bd54, 20:58, dec 4.623 $\rightarrow$ 0.681·0.688), 진단 D1--D3(a4a6257: 계산 한계 아님, 라벨은 학습 가능·일반화 간극, 모호성 아님), 규모 곡선(3478cf2, 23:05) 1k/9.4k/18.7k/37.5k = 0.552/0.682/0.686/0.683 --- 같은 스텝에서 데이터 레버 약하고 포화.
  \item \textbf{그래서(결정):} 다음 레버는 입력(시간 맥락) 쪽(D2).
\end{itemize}

\paragraph{D2. 시간 맥락 $\rightarrow$ 움직임 줄 채택, 2프레임 비디오 불채택}
{\small 09-25 21:42 -- 09-26 02:24 KST}
\begin{itemize}
  \item \textbf{해 봄:} 조사 118def0(21:42): 과거 프레임 1장을 Qwen3-VL 2프레임 비디오로(1순위), 움직임 줄 + 드롭아웃.
  \item \textbf{사용자} (ul 71, 21:48--21:49 KST): \intent{1순위 아이디가 좋다잉}
  \item \textbf{사용자} (ul 71, 21:48--21:49 KST): \intent{저거는 지금 당장 오버리프에 업뎃해줘도 좋겠는데}
  \item \textbf{사용자} (ul 71, 21:48--21:49 KST): \intent{피규어나 이런거에도}
  \item \textbf{해 봄:} E-TC 2$\times$2 사전 등록(13:05:07Z, 0b853c2 22:07): 2프레임 비디오 V $\times$ 움직임 줄 M, se2e, 2,000스텝, 검증 300.
  \item \textbf{결과:} f4f6a49(09-26 00:39): 네 칸 0.683 / 0.696 / 0.708 / 0.721; M 주효과 +0.025 [+0.003, +0.047], 전환 층 +0.027, FULL p95 $-$3.8 \%; V +0.013(< 0.02), FULL p95 +20.4 \%(CPU 전처리).
  \item \textbf{사용자} (ul 80, 00:40 KST): \intent{일단 움직임 줄 채택할게}
  \item \textbf{그래서(결정):} 정본 §83(f00cb59, 00:40): 움직임 줄 \texttt{se2e-\allowbreak{}motion@v1} 채택, 2프레임 비디오 불채택. 발견된 누수: S-E2E 속도가 중앙 차분(미래 0.1 s 포함) $\rightarrow$ 인과 후방 차분 \texttt{se2e\_\allowbreak{}c1}(44c907d, 00:53).
  \item \textbf{결과:} 확인 실험(7dced6a, 02:24): 고친 데이터·시드 2개·검증 1,799에서 합동 +0.032 [+0.022, +0.043], 전이 +0.030 $\rightarrow$ CONFIRMED, §83 유지.
\end{itemize}

\paragraph{D3. MolmoAct 1차 이식 --- 모두 불채택}
{\small 09-26 01:41 -- 10:22 KST}
\begin{itemize}
  \item \textbf{사용자} (ul 83, 01:41--01:42 KST): \intent{VLA도 사진을 보는것도 좋을 것 같구}
  \item \textbf{사용자} (ul 83, 01:41--01:42 KST): \intent{엔드포인트의 움직임을 위에 덧그리는 방식도 뭔가 도움이 될 수 있을 것 같기도 하구}
  \item \textbf{결과:} E-MA1(06355f0 등록) 관문 G0 실패(ee1b000, 02:41): 머리 카메라 FK 투영 명목 100.9/264.5 px, PnP 뒤 86.9/212.1 px(문턱 12/30) $\rightarrow$ 학습 0칸, R2 뒤 재등록.
  \item \textbf{결과:} E-MA1b 로봇 기준 3D 미래 손끝 궤적 보조(ded38d7 등록): 검증 1,799 합동 +0.0031 [$-$0.0022, +0.0085], 전이 0.000 $\rightarrow$ 불채택(2df9bdc, 04:49; §85 보충).
  \item \textbf{결과:} E-CAM3 머리 + 양 손목(시각 토큰 356 $\rightarrow$ 460, 5c46056): +0.0014 [$-$0.0041, +0.0069], FULL 결정 p95 +11.7 \%(84.6 $\rightarrow$ 94.6 ms) $\rightarrow$ 불채택(76ab10c, 07:22); 결정층 입력은 머리 + 활성 손목 유지, 양 손목은 Astra 요청에만(§57 보충).
  \item \textbf{결과:} E-MA3 층별 KV 조건(40895f8): 청크 오차 합동 상대 감소 +2.19 \% [+1.48 \%, +2.88 \%] < 문턱 5 \% $\rightarrow$ 불채택(76ab10c; §84 보충 3).
  \item \textbf{결과:} E-MA2 VLA 입력에 Astra \texttt{edit}(fde4f6a 등록, R2): 돌린 명령 준수율 문장 C1 0.177·화살표 C2 0.001(둘 다 < 0.6); C1은 결정층이 명령 칸을 0.97로 고르지만 청크는 0.28(90$^\circ$)·0.09(180$^\circ$)만 따름(399d1cb, 10:22).
  \item \textbf{그래서(결정):} §84 보충 4: 융합 VLA 입력에 Astra 명령을 넣지 않고 설계 §11 코드 편향만. 결정 토큰이 청크를 조종한다고 가정하지 않는다($\rightarrow$ D5).
\end{itemize}

\paragraph{D4. MolmoAct가 `교훈'으로만 작동 $\rightarrow$ 준비부터, 그리고 실데이터로만}
{\small 09-26 11:25 -- 19:38 KST}
\begin{itemize}
  \item \textbf{사용자} (ul 95, 11:25--11:26 KST): \intent{나는 몰모엑트가 우리가 제대로 준비를 안해놔서 도움이 안되는 교훈ㅇ,로 작동한듯}
  \item \textbf{사용자} (ul 95, 11:25--11:26 KST): \intent{준비를 아ㅖ제대로해서 검증 필요. 시간 걸려두되니까}
  \item \textbf{결과:} R2 준비 관문(6cae25c, 12:04, 학습 없음): 성립 조건 14개; G-cam 표지 중심 0.21 px, G-fk URDF FK 대 시뮬 손끝 9.5 mm 실패, G-trace 현재점 영상 안 0.994, 만들 것 M1--M9.
  \item \textbf{사용자} (ul 99, 14:29 KST): \intent{시뮬거쳤다오는건 절대 반대 몰모엑트 참고해}
  \item \textbf{사용자} (ul 99, 14:29 KST): \intent{추가학습·MolmoAct 이식만 실데이터 (Recommended)}
  \item \textbf{그래서(결정):} 정본 §89(b8c231f, 14:29): 추가 학습·MolmoAct 이식은 실데이터 + 실영상 모델 라벨로, R2 시뮬 MolmoAct 데이터 생성은 중단(파일럿 8편 뒤, 기본 경로 비트 동일 5/5 확인); 조건 준수 연구·폐루프 시험장으로서 시뮬은 유지.
  \item \textbf{결과:} 실데이터 준비(0ca87c5, 15:26): Molmo2-ER 두 팔 포인팅 10,834회, 원시 오류 27--30 \%(대부분 다른 팔); 사전 고정 거르개 v1 실패(유지 0.18) $\rightarrow$ v2 남긴 점 오류 0.037·유지 0.945; 같은 장면·다른 경로 시연은 없음 $\rightarrow$ 조종 데이터(S)는 실물 수집 필요.
  \item \textbf{해 봄:} E-MAR-real 요인 A(7245787, 15:46; 변경 1 3e2f5f6): RB2 실데이터에 MolmoAct 2D 궤적 보조 \texttt{trace5-\allowbreak{}point@v1}.
  \item \textbf{결과:} 825fd37(19:38): 검증 1,799 합동 +0.0047 [$-$0.0006, +0.0103], 전이 +0.0026, RB2 라벨 행 $-$0.003; 보조 궤적 오차 62 px(평균 궤적 기준선 145 px) --- 궤적은 배웠지만 결정 정확도로 옮지 않음, 준수 불변(0.341 대 0.340). 약 3.6 GPU-h.
  \item \textbf{그래서(결정):} §84 보충 11: 불채택. 실데이터 MolmoAct 이식은 대체 경로 시연을 모은 뒤 E-MAR-S로만, 이력 덧그림 O는 보류.
\end{itemize}

\paragraph{D5. 조이스틱 준수 문제 발견 $\rightarrow$ 목표 0.8}
{\small 09-26 10:37 -- 11:22 KST}
\begin{itemize}
  \item \textbf{해 봄:} E-SR0(7b42b97 등록, 학습 없음): 기존 체크포인트에서 전문가에 주는 확정 결정만 반사실로 바꿔 청크 FK 변위가 따르는지(cos > 0.5).
  \item \textbf{결과:} §84 보충 5(10:52 KST): S-E2E 0.340(우연 0.335)·R2 0.431(우연 0.333) $\rightarrow$ WEAK; 참 결정일 때 R2 청크 0.92, 이산 \texttt{phase\_\allowbreak{}id}는 그리퍼 0.95로 따름.
  \item \textbf{사용자} (ul 91, 10:59 KST): \intent{x2점유시작}
  \item \textbf{사용자} (ul 91, 10:59 KST): \intent{준수율이 0.8은 될방법론알 찾아야해}
  \item \textbf{그래서(결정):} 결합은 결정 토큰이 아니라 청크·손끝 기준에 직접 적용(§84 보충 5). E-SR1b 채택 문턱 0.6 $\rightarrow$ 0.8. 방법론 조사(f6c7c3f, 11:22): 1순위 먼 구간 반사실 분기 데이터, 2 권한 가중 구조, 3 CFG w(a), 4 권한 게이트 adaLN; 거리 권한 a(C3)로 모두 재평가.
\end{itemize}

\paragraph{D6. E-SR1b 드롭아웃 + CFG·되짚기 $\rightarrow$ 도달 못 함}
{\small 09-26 11:31 -- 13:51 KST}
\begin{itemize}
  \item \textbf{해 봄:} 사전 등록 ad85725(11:31): R2, E-MA2 C0 레시피에 결정 드롭아웃 0.3 + 결정 CFG(w 1--8)·되짚기 결정 라벨·둘 다, 시드 2개씩 8판.
  \item \textbf{결과:} 12f5902(13:51): 반사실 준수 A\_xy 최고 0.598(되짚기 B, 기준 0.438)·그럴듯한 편집 P\_xy 최고 0.398(기준 0.358); B는 청크 MSE +14.6 \%·그리퍼 일치 $-$0.030으로 비열등 실패; 드롭아웃은 w = 1 준수를 0.362로 낮춤; CFG w 5--8은 청크 MSE +106--357 \%. 약 7.9 GPU-h.
  \item \textbf{그래서(결정):} §84 보충 7 NOT\_REACHED: 드롭아웃 + CFG는 준수 수단 후보에서 뺀다, 결합은 청크·손끝 기준 경유 유지.
\end{itemize}

\paragraph{D7. E-SR1c 먼 구간 반사실 분기 $\rightarrow$ 채택(시뮬)}
{\small 09-26 12:09 -- 14:29 KST}
\begin{itemize}
  \item \textbf{해 봄:} 사전 등록 a79bd64(12:09; 변경 1 2a2ad34, 변경 2 beacc0d): R2 먼 구간(대상까지 $\geq$ 10 cm) 스냅숏에 강제 조이스틱 결정 4개와 그 방향 IK 직선 청크를 짝지은 반사실 분기 128,661행(Isaac 재생 관문 G-br: 추종 중앙 0.9 mm·방향 100 \%·침투 0)을 전문가 배치의 50 \%로; C2는 권한 게이트 adaLN 추가.
  \item \textbf{결과:} e091312(14:29): 먼 구간 반사실 방향 준수 C1 0.991(C0 0.344; A\_z 0.989·$\rho$ 0.992, 분기 없는 bottle\_tray 0.979, 편집 $\pm$30·45$^\circ$ 0.976), 근접 비열등 N1--N5·전 구간 (i)--(iv) 통과, 근접 반사실도 0.868로 따름. 대가: 먼 구간 참 결정 청크 MSE +29 \%·끝점 +4 mm. C2 adaLN 이득 없음($-$0.006)·지연 +9 ms. 약 2.5 GPU-h.
  \item \textbf{그래서(결정):} §84 보충 8 ADOPT\_C1: 먼 구간(a = 1)만 조이스틱 결정 경유 조종 허용, 근접(a = 0)은 VLA 자신의 결정만; §6 투영·§11 편향 $\times$ a; 융합 VLA 레시피에 먼 구간 분기 50 \%. 구현은 Task 17(C5).
\end{itemize}

\paragraph{D8. E-SR1d 같은 레시피를 실데이터로 $\rightarrow$ 도달 못 함}
{\small 09-26 14:59 -- 16:09 KST}
\begin{itemize}
  \item \textbf{해 봄:} 사전 등록 32065de(14:59): S-E2E \texttt{se2e\_\allowbreak{}c1}에 시뮬 없이 URDF IK 기구학 분기 61,993행(관문 G-kin: 거부 11.3 \%, 한계·속도 위반 0), 권한 a = 그리퍼 사건 되짚기 대용(관문 G-a 통과, 참 거리 검증 불가).
  \item \textbf{결과:} 2c93c9b(16:09): 먼 구간 A\_xy 0.339 $\rightarrow$ 0.536(+0.198 [+0.187, +0.209], 시드 0.552·0.521), A\_z 0.519 $\rightarrow$ 0.772, 근접 비열등 통과. 약 2.2 GPU-h.
  \item \textbf{그래서(결정):} §84 보충 10 NOT\_REACHED: 실데이터 체크포인트에는 먼 구간 결정 경유 조종을 옮기지 않는다(청크·손끝 경유 유지); §92 P1(실데이터 준수 $\geq$ 0.8)은 불성립.
\end{itemize}

\paragraph{D9. E-SR1e --- 원인은 적합 부족, 두 팔로 다시}
{\small 09-26 18:59 KST \textasciitilde{} (진행 중)}
\begin{itemize}
  \item \textbf{결과:} 학습 없는 진단(등록 0.1절): E-SR1d C1의 TRAIN 먼 구간 준수 0.549·0.523 = 검증 0.552·0.521 $\rightarrow$ 일반화 차가 아니라 적합 부족; 모든 검증 과제가 학습에 있음, 편별 고르게 `반만' 따름; 2,000스텝 = 0.43 에폭에서 손실이 아직 내려감; 실데이터 전문가는 결정 토큰을 거의 안 씀(C0 참 결정 적중 0.47).
  \item \textbf{해 봄:} 사전 등록 7b076cd(18:59): 팔 A(6,000스텝·분기 50 \%)·D(3,000스텝·분기 75 \%) + L1 대조 결정 안내(E-SR1d C1, w 1.5·2·3, 먼 구간만); 코드 기준 = E-SR1d 등록 커밋(HEAD의 6질문 판은 옛 체크포인트를 못 읽음); RB3는 팔에 넣지 않음(일반화 차 없음).
  \item \textbf{진행 중·대기:} \tentative{결과 문서(\texttt{results/\allowbreak{}sr1e.\allowbreak{}md})는 아직 없다 --- handoff §0의 `진행 중' 줄(메인 GPU 2·3 학습 $\rightarrow$ 평가). 결과에 따라 실데이터 레시피·P1이 바뀐다(ADOPT이면 §84 보충 10 수정).}
\end{itemize}

\paragraph{D10. SigLIP --- 필요 없으면 하지 않는다}
{\small 09-26 10:43 -- 11:03 KST}
\begin{itemize}
  \item \textbf{사용자} (ul 90, 10:43--10:49 KST): \intent{Siglip을 쓸 방법도 찾아봐 어떻게 엮을 수 있을지}
  \item \textbf{사용자} (ul 90, 10:43--10:49 KST): \intent{나는 저걸 우리vla에 쓰는게 어떨가 싶은데}
  \item \textbf{사용자} (ul 90, 10:43--10:49 KST): \intent{그래c,e반영가능성도 검사해보고 결정하자 일단ema2 다되고 나면 그거보고 ㄱ}
  \item \textbf{사용자} (ul 92, 11:03 KST): \intent{만약 시그립이 그래서 필요없다면 굳이 더하란 내말은 무시해두돼}
  \item \textbf{결과:} 조사(4e2e242): 융합 VLA의 영상 탑(Qwen3-VL-4B)이 이미 SigLIP2-Large 계열(24층·폭 1024·패치 16, 동결).
  \item \textbf{그래서(결정):} §84 보충 6(14d78b5, 11:03): 별도 SigLIP 통합은 하지 않음; 손목 패치 (e0)·두 번째 SigLIP 줄 (c)는 준수 문제를 푼 뒤 이득 근거가 있을 때만(기본 보류).
\end{itemize}

\paragraph{D11. 새로움 점수와 결정 확신도 $\rightarrow$ 게이트 없이 기록만}
{\small 09-26 14:36 -- 20:07 KST}
\begin{itemize}
  \item \textbf{해 봄:} E-NOV0(d225a84, 14:36): 백본 특징 k-NN 새로움 점수(학습 없음)로 모델 자신의 큰 오류·새 환경을 미리 아는가.
  \item \textbf{결과:} 7c26faf(15:27): 큰 오류 AUROC 0.570, 90 \% 재현율 호출 감소 13 \%, dr 이동 0.742·과제 빼기 0.751 $\rightarrow$ NONE(흐름 게이트·T\_nov 채널 안 둠, §84 보충 9). 판정 밖: 결정 확신도 0.804(감소 0.450).
  \item \textbf{사용자} (ul 106, 15:37 KST): \intent{결정 확신도는 어쩌다가 발견했음?}
  \item \textbf{사용자} (ul 106, 15:37 KST): \intent{확신도 기반으로 해두면 좋을듯 왜 그런 선택을 했는지도.}
  \item \textbf{사용자} (ul 106, 15:37 KST): \intent{그리고 조금씩 확신도를 능력을 더 높이는 개선도 필요할듯}
  \item \textbf{그래서(결정):} 정본 §95(780ca10, 15:40): 확신도(질문별 1·2등 로그확률 차의 최솟값) 기반 운용 + 선택 이유 기록 + 개선 트랙 --- 결과 뒤 선택이므로 새 표본에서 먼저 검증(E-CONF).
  \item \textbf{해 봄:} E-CONF(6191369, 16:18; 변경 1 f37f8f8·393a184).
  \item \textbf{결과:} cd29570(19:13): R2-새(학습된 적 없는 R2\_TRAIN 실패 편 244편·8,208) AUROC 0.806 [0.796, 0.817]·감소 0.379 통과, S-E2E 실데이터 0.748 [0.724, 0.771]·0.244 실패(문턱 0.75·0.30); 개선 후보 \texttt{temp\_\allowbreak{}joint} +0.041/+0.024(문턱 +0.03 미달).
  \item \textbf{그래서(결정):} §95 보충 1 NONE: 흐름 우선 표시·보수 모드의 트리거로 쓰지 않고 결정마다 `선택 이유' 기록만(훅 \texttt{conf-\allowbreak{}base@v1}, 5a940d2 20:07, 어떤 제어 경로도 읽지 않음).
\end{itemize}

\paragraph{D12. Astra를 추가 학습 교사로(§88) $\rightarrow$ 실데이터 E-AT 설계}
{\small 09-26 14:27 -- 14:46 KST}
\begin{itemize}
  \item \textbf{사용자} (ul 98, 14:27 KST): \intent{아스트라도 추가학습의 일환으로 써야겠음}
  \item \textbf{사용자} (ul 98, 14:27 KST): \intent{본학습은 가정한 지티로 하돼 추가학습으로 실제 반영하는거지}
  \item \textbf{그래서(결정):} 정본 §88(bc923c2, 14:27): 본 학습 = 가정한 정답(규칙 라벨·시뮬 오라클), 추가 학습 = 결합 폐루프에서 Astra가 실제로 낸 판단·수정(DAgger식). 선행 조건: Astra 명령 정확도 측정(E-ACC) --- 틀린 라벨을 배우지 않게.
  \item \textbf{결과:} E-AT 설계(cb81017, 14:46, 문서만): ul 99에 따라 실데이터 전용 --- RB2 실물 시연에 Astra 라벨(구간 품질·J2 진단 위주), 비용 계획 20,000원·상한 25,000원.
  \item \textbf{진행 중·대기:} \tentative{E-AT는 실행 없음: 유료라 사용자 승인 필요(ul 114), \texttt{quality:} 줄은 \texttt{ser-\allowbreak{}A-\allowbreak{}min-\allowbreak{}3}에 없음(결합 계획 G8 --- \texttt{ser-\allowbreak{}A-\allowbreak{}min-\allowbreak{}4}에 묶기).}
\end{itemize}

\paragraph{D13. 직렬화 ser-A-min-3 판 올림 $\rightarrow$ 재학습 대기}
{\small 09-26 15:17 -- 09-27 00:45 KST}
\begin{itemize}
  \item \textbf{해 봄:} 결합 Task 12 한 번의 판 올림: 움직임 줄 + 그리퍼 질문(340c5dc, 15:17), \texttt{prompt\_\allowbreak{}config} 판본·범주 기록과 검사(a761b2b), PH-A1 문구(b931f75), 구간이 그리퍼 사건을 품음·학습 때만 unknown 구간 p 0.3·프롬프트 원천 해시 보강(bce0bd8, 18:45), 재개 시험(fd24d2e, 00:45).
  \item \textbf{결과:} 새 판은 기존 체크포인트를 모두 거부한다 $\rightarrow$ 폐루프·파드 검증이 모두 막힘(C8).
  \item \textbf{진행 중·대기:} \tentative{[결정 필요] ser-A-min-3 재학습(결합 계획 D2 추천: VLA-sim 먼저, VLA-real은 E-SR1e 뒤). book 01 VLA-R3는 [예정].}
\end{itemize}

\paragraph{D14. 우리 VLA는 흔한 VLA가 아니다, 그래도 끝-끝으로 학습}
{\small 09-27 01:4x KST}
\begin{itemize}
  \item \textbf{사용자} (ul 128, 01:4x KST): \intent{나는 VLA 부분은 엔드투 엔드로 학습이 될 필요가 있다고 생각해. 하지만 알겠지만 저 VLA가 우리가 아는 VLA의 역할이 아니라는건 인지하고 있지?}
  \item \textbf{그래서(결정):} 정본 §96 보충 4(a54170f, 01:39): 역할 = 상위 의도(구간 줄) + 영상·상태를 받아 typed 결정(방향·크기·대상·단계 + 쥐기·놓기 시점)을 내는 빠른 조이스틱, 과제를 스스로 해석하지 않음. E2E(E-VLA-E2E, 무료, 새 등록): 결정은 감독된 중간 출력으로 남기고 KI식 전문가 $\rightarrow$ 백본 절연 유지 + 이산 행동 토큰(E1), 정답 $\rightarrow$ 예측 결정 조건 교과(E2), 구간 줄 의도 교란(E4); 결합 고리 안의 VLA DAgger(본학습 자리만).
  \item \textbf{진행 중·대기:} \tentative{선행: ser-A-min-3 단계식 재학습(비교 기준)과 E-SR1e 결과. 행동 토큰은 \texttt{ser-\allowbreak{}A-\allowbreak{}min-\allowbreak{}4} 판 올림이라 E-AT \texttt{quality:} 줄과 한 번에 묶는다. [예정].}
\end{itemize}

% Generated by D:/tools/scratch_qdd/mmrec/render.py from data_e.py - edit the data, not this file.
\subsection{E. 데이터 --- 자체 시뮬이 바닥, 공개 데이터는 픽셀로 바꿔 쓰고, 다양성은 필수}\label{rec:E}
\noindent\textbf{한 줄:} 조종 표본·복구·자기 인식은 공개 대체가 없어 자체 시뮬(참값 모델)이 주원천이다. 남의 로봇 데이터는 행동 라벨이 아니라 픽셀·QA로 바꿔 쓰는 것이 문헌의 방식이고, 지금 과제는 같은 탁자·같은 물체라 다양한 데이터는 필수(OOD 평가 필수). 공개 데이터 변환은 시제품·관문까지, 좌표계 없는 원천은 픽셀 경로(T0 + T5)가 주 추천. 8B 검증 등록(E-OPEN8·E-XEMB8)은 초안.

\paragraph{E1. 자체 시뮬 R2와 소규모 실물 S-E2E}
{\small 09-24 \textasciitilde{} 09-26 09:53 KST}
\begin{itemize}
  \item \textbf{사용자} (ul 66, 09-25 04:18 KST): \intent{직므은 단순이니까 10으로 하되 나중에는 우리가 데이터하고 할떄는 30으로 모을거}
  \item \textbf{그래서(결정):} §62: S-E2E(ROBOTIS 공개 RB1·RB2)는 10 Hz 그대로, 우리 수집 데이터는 30 Hz. §66: R2 생성기 확정(30 Hz·다과제·도메인 무작위화·LeRobot).
  \item \textbf{결과:} R2\_TRAIN(b2b0b3e 08:30, LeRobot 18d4e2f 09:53): 6,000편 생성, 유효 5,126(85.4 \%) --- mug\_tray 99.1 \%·mug\_marker 98.7 \%·bottle\_tray 58.5 \%(실패 874 중 794가 병 놓기); 단계 B 행 1,495,348, LeRobot v2.1 6개 22.13 GB, verify 6/6. S-E2E 42,813행(D1).
  \item \textbf{참고:} 이 데이터가 E-SR1c(시뮬 분기)·E-TEACH-L8(TRAIN 시드)·E-PT의 바탕이다. 지금 시뮬 과제는 탁자 높이 z 0.850·같은 탁자·같은 물체·방해물이다(E6의 출발점).
\end{itemize}

\paragraph{E2. ROBOTIS AI Worker 공개 데이터 전수 조사}
{\small 09-26 15:29 -- 18:43 KST}
\begin{itemize}
  \item \textbf{사용자} (ul 105, 15:29 KST): \intent{1. ai워커에 들어가보면 로보티즈 오픈 데이터 많은데 그중에 쓸만한거 찾아봐 2. 고}
  \item \textbf{해 봄:} 조사(1ef9f2c, 16:20, 학습·유료 0): HF의 FFW/AI Worker 관련 1,150개 데이터셋 메타데이터 전부(ROBOTIS 조직 7개 + 직원 계정 RobotisSW·Dongkkka + 외부), 상위 14개는 프레임 눈 확인. (`2. 고' = E-MAR-real 진행, D4.)
  \item \textbf{결과:} 1순위 \texttt{Dongkkka/\allowbreak{}ffw\_\allowbreak{}bg2\_\allowbreak{}rev4\_\allowbreak{}pickup\_\allowbreak{}obj\_\allowbreak{}1127\_\allowbreak{}total2}(806편, 우리 기종 BG2, 카드 apache-2.0; 같은 공구 벽에서 지시문이 대상을 고르는 쌍 601·같은 물체를 다른 팔로 457), 2순위 RobotisSW 과제 카탈로그(322과제·7,364편, apache-2.0, 지시문은 번호뿐), 3순위 Merged\_PickPlace 병/캔(두 대상, 라이선스 없음). 같은 팔·같은 장면의 대체 경로 시연은 어디에도 없음; 실패 데이터는 외부 사용자 것뿐.
  \item \textbf{결과:} 2차(66f2834, 18:43, 사용자 후속 질문 `VLA 오픈 데이터'): 공식 AI Worker VLA 데이터 = HF ROBOTIS 7개(+ GR00T 미세조정 모델 6개). 약 10,500편 FK 채굴: 같은 시작·같은 목표에서 $\geq$ 5 cm·$\geq$ 40 px 갈라지는 쌍 RB2 420·bg2 192·pickup\_obj 약 900; $\geq$ 10 cm·$\geq$ 80 px는 라이선스 있는 원천에서 약 100쌍뿐. 프레임상 의도적 우회가 아니라 사람의 좌우 흔들림.
  \item \textbf{그래서(결정):} 이 쌍들은 오프라인 궤적 추종 평가(n $\geq$ 300)·두 대상 평가(pickup\_obj 601 + RB2 426)에만 쓰고 MolmoAct 조종 학습 데이터로는 쓰지 않는다. 1순위 데이터는 RB3로 추가(E3).
\end{itemize}

\paragraph{E3. RB3 추가 --- 새 데이터 판 se2e\_c2}
{\small 09-26 18:52 -- 18:56 KST}
\begin{itemize}
  \item \textbf{결과:} 2f973ba(18:56): RB3 806편·104,470프레임·3.31 GB(verify 오류 0) $\rightarrow$ 변환 627편·16,721행(그리퍼를 안 쓴 빈 편 179 제외, 프레임 확인), train 15,906 / val 815. RB1·RB2와 해시 일치 0·궤적 최근접 0.351 rad. 조사와 다른 사실: 손목 영상 회전 불필요, 파켓 task\_index 58편 오류(episodes.jsonl 사용), \texttt{release\_\allowbreak{}visible} 240/627(눈 대조 24/24).
  \item \textbf{그래서(결정):} \texttt{se2e\_\allowbreak{}c2} = \texttt{se2e\_\allowbreak{}c1} + RB3(RB1·RB2 행은 바이트 동일). 카드는 apache-2.0이지만 직원 개인 계정이라 내부용 표시. E-SR1e는 일반화 차가 없어 RB3를 팔에 넣지 않았다(D9).
\end{itemize}

\paragraph{E4. 상위 VLM 학습 데이터 조사 --- 시뮬이면 더 좋다}
{\small 09-26 21:0x -- 21:17 KST}
\begin{itemize}
  \item \textbf{사용자} (ul 111, 21:0x KST): \intent{저거에 맞는 데이터를 찾는 작업을 꼼꼼하게 해보자 심이면 더 좋을 것 같구}
  \item \textbf{결과:} 조사(73787c6, 21:17): 약 40개 데이터셋 확인, 7종 파드 표본 확인. 판정: 조종 표본·복구(공개 수치 교정 데이터 없음)·자기 인식(AI Worker 공개 시뮬 데이터 0건)은 자체 Isaac 시뮬 + 참값 모델이 주, 공개는 장면 다양성·일반 그라운딩 보조. 상위 5: MolmoBot-data(RBY1)·BEHAVIOR-1K 2025·RoboTwin 2.0·RefSpatial·InternData-M1(게이트·NC). 막힘: 게이트 동의, NC·ND 라이선스는 Astra 전송 보류.
\end{itemize}

\paragraph{E5. 남들은 맞지 않는 로봇 데이터를 어떻게 쓰나 --- 35B 데이터 양}
{\small 09-26 22:2x -- 22:29 KST}
\begin{itemize}
  \item \textbf{사용자} (ul 115, 22:2x KST): \intent{기존 데이터를 쓸 방법을 찾아봐 다른 연구들에서는 우리와 같지는 않겠지만, 다른사람들은 자기에 맞지 않은ㄴ 로봇데이터를 어떤식으로 튜닝 혹은 활용을 했는지 찾아봐 35b학습시키는데 데이터가 충분해야하다고 생각하거든 거기다 범용성도 있으려면 말이야}
  \item \textbf{결과:} 조사(2373a02, 22:29): 남들은 맞지 않는 로봇 데이터를 행동 라벨로 쓰지 않고 (a) 기체 무관 영상 표현(MolmoAct 그리퍼 2D 궤적 + 깊이 토큰, Embodied-R1 점, HAMSTER[기초]) 또는 (b) QA·추론 재라벨(Robo2VLM, InternVLA-M1, RoboBrain 2.0)로 바꾸고, 출처 표시(X-VLA)·일반 데이터 5--25 \% 공동학습으로 범용성을 지킨다. Vlaser(ICLR 2026): 남의 영역 데이터만으로는 조종 이득 없음(41.8 $\rightarrow$ 43.2 \%, 같은 영역 65.1 \%).
  \item \textbf{그래서(결정):} 자체 시뮬 주 판정 유지, 공개는 픽셀·카메라 좌표 인식 QA로(출처·좌표계·내부값을 글로). 35B 1차 약 40만 표본(자체 조종 25 \%·자체 인식 25 \%·공개 인식 25 \%·추론 10 \%·일반 공간 10 \%·일반 5 \%) [추정]. 무료 검증 E-XEMB8은 설계만, L8 결과 뒤.
\end{itemize}

\paragraph{E6. 다양한 데이터는 필수 + OOD 평가}
{\small 09-26 23:1x -- 23:36 KST}
\begin{itemize}
  \item \textbf{사용자} (ul 116, 23:1x--23:2x KST): \intent{근데 쓸만한 거는 가공해서 조종명령 답하게 바꾸는건 더 좋은건 아니야?}
  \item \textbf{사용자} (ul 116, 23:1x--23:2x KST): \intent{왜 그러냐면 우리의 테스크는 동인 높이 책상에 동일한 장애물들에 다 동일하잖아 무조건 다른데이터가 필수로 있기는 해야해}
  \item \textbf{그래서(결정):} 정본 §97 보충 1(94c0440, 23:36): 상위 VLM 학습에서 다양한 데이터는 필수(넣을지를 실험으로 정하지 않음, E-XEMB8은 형태·비율만). (a) 공개 변환 격상 --- 인식 QA + 좌표계 명시 조종 C$'$(\texttt{source}·\texttt{frame: base\_\allowbreak{}<기체>}, 그 기체 자기 정보 블록 자동 생성, 그리퍼 사건에서 우리 단계 어휘로 자름); (b) 자체 생성기 다양화 --- 탁자 높이(목표 $\pm$10--15 cm)·크기·색, MolmoSpaces 물체, 방해물 0--6·혼동체, 받침, 조명·질감, 카메라 장착 흔들림(머리 각도는 고정), 새 과제 변형 4; 기본 경로는 바이트 동일.
  \item \textbf{그래서(결정):} OOD 평가(학습 절대 금지): 보호 분할 \texttt{ood} 시드 70000--70999 --- OOD-H 높이·OOD-O 물체 범주 20 \% 보류·OOD-D test 풀 방해물·OOD-X 공개 기체 보류; 지표는 L8과 같고 프레임 누수율 0이 관문.
  \item \textbf{진행 중·대기:} \tentative{사용자 확인 필요로 남은 것: 탁자 높이 범위를 리프트(\texttt{set\_\allowbreak{}lift})로 맞춰도 되나, 게이트 데이터 HF 동의 여부. 생성기 다양화·OOD 세트는 [예정].}
\end{itemize}

\paragraph{E7. 공개 데이터 변환 방법 --- MolmoAct 방식만으로 되나}
{\small 09-27 00:14 KST (커밋 시각)}
\begin{itemize}
  \item \textbf{사용자} (ul 119, 00:3x KST): \intent{1번을 할 방법을 더 자세하기 찾아봐야해}
  \item \textbf{사용자} (ul 119, 00:3x KST): \intent{몰모엑트가 한방식데로 데이터를 하면 안되나?}
  \item \textbf{사용자} (ul 119, 00:3x KST): \intent{우리는 몰모엑트방식으로는 데이터가 부족한가?}
  \item \textbf{결과:} \texttt{public\_\allowbreak{}data\_\allowbreak{}conversion\_\allowbreak{}howto} + CPU 시제품 \texttt{tools/\allowbreak{}xemb/\allowbreak{}}(624fdd6, 00:14; 시험 31개, 유료 0·GPU 0): MolmoAct 공개분(CC BY, 약 1,000만 프레임)은 픽셀 점·궤적은 넘치지만 기저 좌표 조종(C$'$)·metric 3D·평면은 0이라 그것만으로는 부족; 보정 있는 원천은 포인팅 대신 FK + 보정 투영. C$'$ 주원천 MolmoBot Franka(위에서 잡기 93 \%) $\rightarrow$ RoboTwin(73 \%) $\rightarrow$ BEHAVIOR(29 \%). 시제품 관문: EE가 그리퍼 위 BEHAVIOR 37/37·MolmoBot 자동 90.2 \%(눈 $\geq$ 97 \%)·RoboTwin 눈 9/9, 단계 $\leftrightarrow$ 계획기 94.2 \%, 탁자 평면 61/61, 분할 $\cap$ 깊이 3D 중심 9.1 cm 불통.
  \item \textbf{그래서(결정):} 3D 답은 참값(GT)에서만 만든다. 변환 결과 P 1,145·C$'$ 54(시제품). 세트 구성은 E-OPEN8·E-XEMB8 초안으로(E9).
  \item \textbf{참고:} user-log 119의 시각 표기는 15:3x UTC(= 09-27 00:3x KST)이고 커밋은 00:14 KST라 순서가 맞지 않는다 --- 기록 그대로 둔다.
\end{itemize}

\paragraph{E8. 좌표계 없는 공개 데이터 --- 가장 효율적인 방법 채택}
{\small 09-27 00:4x -- 00:49 KST}
\begin{itemize}
  \item \textbf{사용자} (ul 124, 00:4x KST): \intent{제일 효율적인 방법으로 채택을 해야해 그것도 연구 자료 뒤져서 저 세가지 방법외에 하나 더 찾아봐 2개도 괜찮고}
  \item \textbf{사용자} (ul 124, 00:4x KST): \intent{너가 아까 말한 좌표계에서 픽셀로만 내는것도 방법이 될 수는있고}
  \item \textbf{사용자} (ul 124, 00:4x KST): \intent{근데 저 4개의 방법론은 오픈소스데이터로 일단 해야해 이해했으?}
  \item \textbf{사실 확인:} 기존 트랙(다른 에이전트, \texttt{tools/\allowbreak{}xemb/\allowbreak{}routes.\allowbreak{}py}): T0 픽셀만, T1 좌표계 만들기(FK 3D + 검출 2D $\rightarrow$ PnP), T2 암묵 학습, T3 단안 metric 깊이.
  \item \textbf{결과:} 조사(a455fc8, 00:49, 코드·GPU·유료 0): 새 방법 T4 로봇 몸 정합 보정(CtRNet-X ICRA 2025 --- DROID 공식 2025 재보정 약 3.6만 장면에 쓰임, 렌더 IoU 0.019 $\rightarrow$ 0.836), T5 점 추적 픽셀 감독(CoTracker3 ICCV 2025, Magma CVPR 2025가 OXE 약 97만 궤적 라벨에 사용) + 관절값 그리퍼 사건으로 집기·놓기 점. 스테레오·다시점 재구성은 T3 변형, 잠재 행동·Cloak은 뺌.
  \item \textbf{그래서(결정):} 정본 §97 보충 3(추천안, 검증 후 확정): 보정 없는 공개 데이터(AI Worker·OXE)에 먼저; 주 = 픽셀 경로 T0 + T5(E-PT \texttt{point} 형식 C-pt 포함, m 값 없음), 보조 = T1(+ T4, 재투영 $\leq$ 5 px 관문 통과분만 투영·C$'$), T3는 관문 전 서열 QA만, T2는 방법이 아니라 결과(ND-1로 판정). 공유 앞단은 추적(T5).
  \item \textbf{진행 중·대기:} \tentative{확정 절차: L단계(라벨 정확도 --- 보정 있는 공개 데이터는 보정을 가린 참값, AI Worker·OXE는 클릭 약 1,000점) $\rightarrow$ T단계 8B 팔 A·PX·CAL(·D3), E-PT와 같은 DEV·OOD-H + OOD-X-open. 둘 다 아직 실행 전.}
\end{itemize}

\paragraph{E9. 8B 검증 등록 초안 E-OPEN8·E-XEMB8}
{\small 09-27 KST (미커밋 초안)}
\begin{itemize}
  \item \textbf{진행 중·대기:} \tentative{\texttt{docs/\allowbreak{}stage3/\allowbreak{}prereg\_\allowbreak{}open8.\allowbreak{}md}·\texttt{prereg\_\allowbreak{}xemb8.\allowbreak{}md}는 작업 트리에만 있는 초안(`등록 아님, 실행 안 함')이고 커밋되지 않았다.}
  \item \textbf{참고:} E-OPEN8 초안: O-A 공개 데이터만으로 LoRA한 8B가 우리 요청 형식에서 대상에 다가가나(영샷 대비), O-B L8 자체 데이터 + 공개 세트를 같은 걸음 수로 섞으면 DEV·OOD-H가 나아지나; 좌표 누수 0; E-PT와 같은 모델·816걸음·같은 스냅숏·같은 채점기; 팔당 약 1.4 GPU-h [추정].
  \item \textbf{참고:} E-XEMB8 초안(§97 보충 1 개정판): 같은 상태·같은 계산량에서 L8 학습 행 일부를 (B) 공개 인식 QA, (D) 공개 C$'$, (E) 둘 다로 바꿔 가르기; 선행 = E-PT 결과(인터페이스)·관문 통과 세트. 3팔 약 4.2 GPU-h [추정].
\end{itemize}

\paragraph{E10. 좌표계 없는 오픈소스 --- 세 트랙 구현, 오픈소스 먼저, 점 추적 공용}
{\small 09-27 01:0x -- 02:5x KST}
\begin{itemize}
  \item \textbf{사용자} (ul 123, 01:0x -- 02:5x KST): \intent{아니 아까부터 계속 이얘기 했잖아 좌표계 없는거는 어떻게 할건지 오픈소스를 어떻게 할건지 해야한다니까?}
  \item \textbf{사용자} (ul 123, 01:0x -- 02:5x KST): \intent{좌표계가 없는걸 좌표계를 만들건지 / 그냥 이미지를 엄청 ㅁ낳이 핛브햇거 좌표계를 뭔가 너머에서 알 수 있게 하던지 / 뎁스를 근사화하는 비슷한 방식으로 해서 하던지}
  \item \textbf{사용자} (ul 123, 01:0x -- 02:5x KST): \intent{일단 오픈소스를 가지고 학습을 할건데 어떤걸 써서 어떻게 할거임? 대규모는 아니고 소규모로 해두돼}
  \item \textbf{사용자} (ul 123, 01:0x -- 02:5x KST): \intent{점추적은 T,1,3에 다 들어갈 수 있을 것 같은데 아닌가?}
  \item \textbf{결과:} \texttt{tools/\allowbreak{}xemb/\allowbreak{}} + howto 9--12절(7faca71): RB2·MolmoAct OXE에 먼저, 보정 있는 공개 세트는 보정 숨긴 정답 대조만. 공용 점 추적 앞단(CoTracker3). T1: 정확한 2D면 0.81 px, 포인팅·추적 2D면 24--26 px(불통). T3(MoGe-2): 물체 17.5 cm·평면 7.7 cm(불통). E-OPEN8 오픈만 6,500행.
\end{itemize}

\paragraph{E11. 작전 T --- 깊이 추정은 버리고 두 안(T1+T4 대 C$'$)을 끝까지 키운 뒤 결정}
{\small 09-27 03:0x KST \textasciitilde{} (최종 결정 12:53 KST)}
\begin{itemize}
  \item \textbf{사용자} (ul 144, 03:0x KST): \intent{일단 내가 시간을 더 줄테니까 계속 도전해보다가 10시간 줄게 안되면 잘되는 탑2에 다 몰아주고 그다음에 탑1에 다 몰아주는걸로 할게}
  \item \textbf{사용자} (ul 144, 03:0x KST): \intent{아니다 깊이 추정은 그냥 아닌게 맞는 것 같아 버리고 2개중에 고민하고 계속 각각을 최대로 발전시킨 후에 10시간 뒤에 최종 결정하자. 이걸 작전명 T라고 할게}
  \item \textbf{결과:} T4 실루엣 정렬(d2b2251·caea370): RB2 보류 경계 2.32 px, 프레임 96 \% $\leq$ 5 px $\rightarrow$ RB2 투영 점 18,636. 명목 카메라도 4.65 px라 T1의 21--42 px는 포인팅 잡음. RB3 2.87 px(73 \%), RB1 5.48 px(44 \%, 팔 마스크 병목)(5033591). C$'$ 물체 이름을 과제 문장에서(`the object' 누수 제거, 746396c). T3 재검토: MoGe-2 + FK 척도 정렬도 물체 19 cm·평면 12 cm $\rightarrow$ 폐기 유지. E-OPEN8 O-A 옛 시뮬 예비: DEV 106.8 mm(영샷 0.85배, NONE), 누수 1/305.
  \item \textbf{진행 중·대기:} \tentative{최종 결정 09-27 12:53 KST. 선택은 L8-X에서 같은 조건 짝 비교 뒤(ul 143). 논문 본문 3.3절·5절에 `결정 대기'로 반영.}
\end{itemize}

\paragraph{E12. E-DIST8 --- 데이터 3수준 $\times$ 공개 픽셀 유무 6팔}
{\small 09-27 03:3x KST}
\begin{itemize}
  \item \textbf{사용자} (ul 141, 03:3x KST): \intent{A·B·C·D 모두에 픽셀을 넣으면 픽셀이 도움인지 방해인지 확인하는 것까지 넣는걸로 하자}
  \item \textbf{그래서(결정):} E-DIST8(8B, 무료, 같은 걸음·표본 수): A(L8)·A+px, B(L8-D)·B+px, C(L8-D + 공개 3D: BEHAVIOR 계량 3D QA + C$'$)·C+px. 평가 DEV·OOD-H(좁게·넓게)·OOD-O/D. 픽셀 효과는 세 수준 각각 짝 비교. L8-D 준비 뒤 등록. 논문 6절 `이번 마일스톤'에 반영.
\end{itemize}

\paragraph{E13. L8-X 환경 모음 --- 시뮬을 약 10배 다양하게, 데이터셋은 1.5년 규칙 해제}
{\small 09-27 03:5x -- 04:5x KST}
\begin{itemize}
  \item \textbf{결과:} L8-D 생성기 등록(e4dced8): 한 Isaac 프로세스에 한 높이, 높이별 작업 상자 = 머리 시야 띠 $\times$ IK 도달 띠, 관문 G-H --- 고정 리프트 0.74--0.98 m 통과. L8-X 설계(dfeb965): 받침면 6(탁자·조리대·계단식 선반·낮은 탁자·바닥 통/상자·윗판 수납장 = 한계 탐침), 과제 약 12, 방해물 0--8, 외관 확장. 물체·과제(e6b2e1e): 받침대·파란 머그·작은 빨간 컵·회색 통, 학습 X 과제 7종. 보류 분할 학습 전 고정(d37634b): OOD-O 작은 컵, OOD-T 보류 조합 2종, OOD-S는 가구 장면 들어올 때. 여러 단계 과제 2종(aeb37f3).
  \item \textbf{사용자} (ul 147, 04:5x KST): \intent{데이터는 1년반해제해줄게}
  \item \textbf{그래서(결정):} CLAUDE.md 예외 줄(99a98d9·d0fdc10): 데이터셋에는 1.5년 규칙을 적용하지 않음(DROID·RH20T·OXE 등 사용 가능), 논문·방법 근거에는 유지; 라이선스·게이트·신뢰도·Astra 전송 제한은 그대로. 논문 3.3절에 한 문장.
\end{itemize}

% Generated by D:/tools/scratch_qdd/mmrec/render.py from data_f.py - edit the data, not this file.
\subsection{F. 과정과 규칙 --- 기록, 자체 검사, 영상, 절약, 논문 갱신, 이름}\label{rec:F}
\noindent\textbf{한 줄:} 모든 결정과 실수를 MD·기록책에 남기고(함정 P01--P123), 모든 실험은 사전 등록 + 시작 전·도중·끝 자체 검사, 폐루프 판은 영상 필수, 유료는 명시 승인 뒤. 논문은 가설 기반 완성판 + 빨간 줄 마인드맵 두 판을 4시간마다 갱신. 논문 이름은 PaceNotes(저장소·코드 이름 harvest 유지).

\paragraph{F1. 모든 것을 MD로 --- 같은 실수를 반복하지 않게}
{\small 09-23 \textasciitilde{} 09-24 04:58 KST}
\begin{itemize}
  \item \textbf{사용자} (ul 12·16, 09-23 첫 세션): \intent{초안 작성 버전도 MD로 다 기록한다(이때는 이렇게 했었다). 같은 실수를 반복하지 않기 위해서다.}
  \item \textbf{사용자} (ul 12·16, 09-23 첫 세션): \intent{사용자가 의도하고 말한 것은 빨간색, 나머지는 검정색.}
  \item \textbf{사용자} (ul 20, 09-24 04:50 KST 경): \intent{꼭 MD를 쓰면서 진행한다. MD에 기록해서 같은 실수를 반복하지 않는다.}
  \item \textbf{그래서(결정):} CLAUDE.md(작업 규칙)·\texttt{docs/\allowbreak{}user-\allowbreak{}log.\allowbreak{}md}(사용자 발언 전부)·\texttt{docs/\allowbreak{}draft-\allowbreak{}log.\allowbreak{}md}(판마다 한 일·틀린 것·교훈)·정본 \texttt{docs/\allowbreak{}design/\allowbreak{}00-\allowbreak{}interfaces.\allowbreak{}md}(뒤 절이 앞 절보다 우선). 마인드맵의 빨간 줄은 user-log 글자와 일치해야 하고 \texttt{tools/\allowbreak{}intent\_\allowbreak{}check.\allowbreak{}py}로 확인한다.
\end{itemize}

\paragraph{F2. 논문은 두 판 --- 논문형 본문과 빨간 줄 마인드맵}
{\small 09-24 04:58 KST \textasciitilde{} 09-27}
\begin{itemize}
  \item \textbf{사용자} (ul 21, 09-24 04:58 KST): \intent{Overleaf는 절대 논문이 아니다. 마인드맵 그리듯 초안을 잡는 느낌이다.}
  \item \textbf{사용자} (ul 26, 09-24 12:32 KST): \intent{내용이 얼마나 많고 적음은 따지지 말고 지금 한 것까지만 해서 글을 좀 다듬어 줘봐 지금 있는 내용에서만 글을 다듬어 달라는 얘기야 막 어, 추가로 뭘 넣거나 이런다는 얘기가 아니라}
  \item \textbf{사용자} (ul 26, 09-24 12:32 KST): \intent{실제 논문처럼 글을 써주면 돼.}
  \item \textbf{그래서(결정):} 본문 \texttt{paper/\allowbreak{}main.\allowbreak{}tex} + \texttt{sec/\allowbreak{}*.\allowbreak{}tex} = 논문형 한국어 글(있는 내용만), 마인드맵 \texttt{paper/\allowbreak{}mindmap.\allowbreak{}tex} + \texttt{mindmap/\allowbreak{}sec/\allowbreak{}*.\allowbreak{}tex} = 빨간 줄(사용자 의도)을 보관하며 계속 갱신.
  \item \textbf{사용자} (ul 63, 09-25 03:18 KST): \intent{지금까지 된 거 어떤식으로 할건지 해서 실제 논문형식으로 오버리프 다시 확 다듬어줘 아직 확실한 자료가 아닌건 임시로 넣어두고 하는 것도 괜찮아}
  \item \textbf{사용자} (ul 88, 09-26 03:02 KST): \intent{지금의 가설을 기반으로 지금 아직 진행안된거 기바ㅏㄴ으로도 해서 제일 이상적인걸로 정리최대해서 논문 ㅈㅇ리해서 지금바로 줘봐보바 당연히 나중에 계속 업뎃 동일하게해주고}
  \item \textbf{그래서(결정):} 정본 §85·CLAUDE.md(13f3495, 03:02): 논문은 가설 기반 이상적 완성판 --- 안 한 실험은 가설·등록 계획·예상으로 `예정' 표시(주황), 숫자 날조 금지. 첫 판 745af38(03:19).
  \item \textbf{참고:} ul 113(B5)으로 확정 전 가설도 그 자리에서 논문(주황)·마인드맵에 올린다. 
  \item \textbf{사용자} (ul 132, 09-27 02:0x KST (통제자 경유)): \intent{자 논문을 우리 2개 놓고 있잖아? 그 한개는 내가말한말 있는거에는 엄청다 적는걸로 하자 이렇게 해봤는데 ㅇ래서 이렇게 하기로 했다 이런식의 상세한것까지. 단 원본 논문은 이제 진짜 다듬는걸 한번 실행을 해보자ㅣ}
  \item \textbf{사용자} (ul 133, 09-27 02:0x KST (통제자 경유)): \intent{그 피규어들도 전부 다 돌면서 새로 만드는거나 삭제, 수정 다 해줘봐봐 메인 피규어도 마찬가지}
  \item \textbf{그래서(결정):} 두 판을 나눔: 본문은 다듬기 판(14930f7 --- 이야기 한 줄기, 실험 장부성 내용은 부록·마인드맵으로, 50 $\rightarrow$ 14쪽)과 그림 정리(논문 에이전트), 이 마인드맵은 상세 결정 이력(5절, 17395b3)이고 제목을 `PaceNotes --- 연구 기록·마인드맵 (상세 결정 이력)'으로.
\end{itemize}

\paragraph{F3. 논문 갱신 주기: 1시간 $\rightarrow$ 4시간, 큰 결과는 수시}
{\small 09-25 18:25 KST \textasciitilde{}}
\begin{itemize}
  \item \textbf{사용자} (ul 68·69, 09-25 18:25--18:28 KST): \intent{지속적으로 오버리프 업데이트 하면서 하는중?}
  \item \textbf{사용자} (ul 68·69, 09-25 18:25--18:28 KST): \intent{1시간마다 논문 업데이트하기로해 앞으로는 꼭}
  \item \textbf{사용자} (ul 73, 09-25 23:04 KST): \intent{논문 업뎃주기는 4시간으로 하자}
  \item \textbf{그래서(결정):} CLAUDE.md(93c68d7, 585786d): 4시간 주기(매 4시간 :17), 큰 판정은 주기와 별도 반영. 정기 갱신 예: cd5210f(09-26 01:07), 54b9205(08:57), 5ea8585(12:56), 858ac6c(18:50), 9d12a4a(20:52), 1d3c86c(09-27 00:51); 수시 갱신 054f250(H-L), 35c9fbb(L8).
\end{itemize}

\paragraph{F4. 조사 규칙과 보고 시각}
{\small 09-23 \textasciitilde{} 09-25 23:03 KST}
\begin{itemize}
  \item \textbf{사용자} (ul 14, 09-23): \intent{신뢰도 낮은 논문과 깃 저장소는 최대한 쓰지 않는다. 쓰느니만 못하다.}
  \item \textbf{사용자} (ul 70, 09-25 19:03 KST): \intent{문제가 있다면 은 만약에 성능이 너무 낮다면 은 기존에 최근 1년 반 이후 논문에서는 그러한 성능들을 비슷한 논문을 찾아가지고 그러한 성능들을 어떻게 올렸는지 찾아봐서 적응을 해도}
  \item \textbf{그래서(결정):} 문헌은 1년 반 이내(큰 주제의 기초 문헌은 `기간 밖, 기초' 표시), 학회·인용·별로 신뢰도 확인. 성능이 낮으면 선행 연구의 개선법을 새 판·새 사전 등록으로 적용(문턱을 결과 뒤에 바꾸지 않음) --- 예: 준수율 0.8 조사(D5), E-SR1e의 L1 대조 안내(D9).
  \item \textbf{사용자} (ul 72, 09-25 23:03 KST): \intent{한국시간으로 해줘 앞으론}
  \item \textbf{그래서(결정):} 사용자 보고는 KST, 저장소 기록은 UTC(d734d00). 이 5절은 읽는 사람을 위해 KST로 적었다.
\end{itemize}

\paragraph{F5. 실험마다 자체 검사}
{\small 09-26 02:41 KST \textasciitilde{}}
\begin{itemize}
  \item \textbf{사용자} (ul 87, 02:45 KST): \intent{매번 자체검사를 해야해 알겠지?}
  \item \textbf{그래서(결정):} CLAUDE.md(ed478e4): 모든 실험(유료·GPU)은 (1) 시작 전 --- 바꾸는 설계 결정, 이 표본으로 가를 수 있나, 더 싼 사전 실행을 사전 등록에; (2) 도중 --- 관문 결함(버그·무효 > 10 \%·교란·비용이나 시간 2배·관문 실패)이면 멈추고 재설계 $\rightarrow$ 등록 변경을 커밋한 뒤 재개; (3) 끝 --- 검증 명령으로 확인 뒤 보고.
  \item \textbf{결과:} 실제 적용 예: 탐침 재범위(A3), E-ACC 변경 1--6(A6·A7), Astra 단독 변경 1--3(A8·A9), E-SR1c 변경 1·2(D7), E-PT 변경 1(B7), E-STRIP8 변경 1(B9). 규칙 이탈도 결과 문서에 적는다 --- heredoc·\texttt{python -\allowbreak{}}·\texttt{python -\allowbreak{}c}를 쓴 일이 여러 에이전트에서 되풀이됨(함정 P03).
  \item \textbf{참고:} 유료 쪽 `승인 불필요' 부분은 ul 114가 대체했다(A10). 자체 검사 의무는 그대로.
\end{itemize}

\paragraph{F6. 결과·실수를 책처럼 --- 연구 기록책}
{\small 09-26 03:34 -- 03:59 KST}
\begin{itemize}
  \item \textbf{사용자} (ul 89, 03:34 KST): \intent{결과들은 모두 효율적으로 정보 저장해둬 나중에 같은씰수안하게 동일 md를 써도되고 효율적으로 분산해서 책처럼 기록 남겨둬두되고}
  \item \textbf{사용자} (ul 89, 03:34 KST): \intent{업데이트 기록 필수임}
  \item \textbf{그래서(결정):} \texttt{docs/\allowbreak{}book/\allowbreak{}}(e15270a, 03:40): 01 실험 장부·02 함정과 예방 규칙·03 결정 연표·04 데이터/코드/파드 지도·05 비용 장부; 06 갱신 기록 필수(4117851, 03:59). 결과·실수가 생기면 같은 커밋에 책도 갱신, 새 작업 전에 함정 장을 먼저 읽는다.
  \item \textbf{결과:} 지금 함정 장은 P123까지 있다. 예: P80(경로 없이 커밋해 남의 파일을 쓸어 담음 $\rightarrow$ 늘 경로 지정 커밋), P101(불리한 기준선), P108(크레딧 소진을 빈 답으로 셈), P113(외삽 예상 상태), P122(IK 잠김 막힘 행), P123(vLLM 환경·병합 샤드 색인).
\end{itemize}

\paragraph{F7. 영상 규칙}
{\small 09-26 19:21 KST \textasciitilde{}}
\begin{itemize}
  \item \textbf{사용자} (ul 131, 09-26 19:2x KST (통제자 경유)): \intent{헐 그거. 영상들도. 무조건 다 데이타에 저장해둬}
  \item \textbf{사용자} (ul 131, 09-26 19:2x KST (통제자 경유)): \intent{필수야 심서 아스트라가 조금이라도 돌릴때 영상 저장.}
  \item \textbf{사용자} (ul 131, 09-26 19:2x KST (통제자 경유)): \intent{노노 그냥 아까 아스트라 단독만 일단}
  \item \textbf{그래서(결정):} Astra 단독 실제 편 7개 영상(편마다 머리 | 오른손목 연속 영상 + 모델 시점 영상)을 파드 \texttt{/\allowbreak{}data/\allowbreak{}harvest/\allowbreak{}videos/\allowbreak{}}에 저장(4d54506). 이후 폐루프 시험의 등록·결과는 영상을 필수로 적는다: E-OpenVLM 대리(\texttt{videos/\allowbreak{}open\_\allowbreak{}vlm\_\allowbreak{}proxy/\allowbreak{}}), E-TEACH-L8(\texttt{videos/\allowbreak{}teach\_\allowbreak{}l8/\allowbreak{}}), E-PT·E-STRIP8 등록(폐루프 영상).
\end{itemize}

\paragraph{F8. 절약과 작업 위치}
{\small 09-24 17:30 KST \textasciitilde{}}
\begin{itemize}
  \item \textbf{사용자} (ul 43, 09-24 17:30 KST): \intent{넣는 건데 2번, 3번, 3번은 일단 비용 한도 없이 그냥 할 거니까 진행해 주되 최대한 비용 문제 없게 진행을 해주고 2번 경우에는 어, 사용해도 돼 1번의 경우에는 일단 후순위로 미루자 내가 나중에 줄게}
  \item \textbf{사용자} (ul 45·51, 09-24 18:24 · 21:20 KST): \intent{H200 서버 사용할 때 무조건 데이터 경로에다가 사용해줘야 돼.}
  \item \textbf{사용자} (ul 45·51, 09-24 18:24 · 21:20 KST): \intent{필요한 큰거나 기본은 전부다 data 경로인건 알지? 쿠베의 h200의}
  \item \textbf{사용자} (ul 45·51, 09-24 18:24 · 21:20 KST): \intent{폴더 이름은 harvest 여야지}
  \item \textbf{그래서(결정):} 유료 한도 약 10만 원(ul 76) $\rightarrow$ 실험별 상한·80 \% 정지 $\rightarrow$ 이제 명시 승인 뒤에만(ul 114, A10). 무료 모델로 먼저 배선·버그를 잡고 유료는 가를 수 있는 표본만(탐침 재범위, Astra 단독 단계식). 파드 파일은 모두 \texttt{/\allowbreak{}data/\allowbreak{}harvest/\allowbreak{}} 아래, 로컬 임시 파일은 \texttt{D:\textbackslash{}tools\textbackslash{}scratch\_\allowbreak{}qdd}(C 드라이브 금지).
\end{itemize}

\paragraph{F9. 논문 이름 PaceNotes}
{\small 09-27 01:46 KST (커밋)}
\begin{itemize}
  \item \textbf{사용자} (ul 129, 09-27 01:5x KST): \intent{우리 이거 일단 논문을 정리를 할건데 이 논문의 이름을 뭐라할까? 내가 이논문이름을 원래 harvest로 했던 이유는 seed라고 그래스프잰뿌려리고 그리고 mill 그중에 장애물이나 잡기쉬운걸로 하고 하베스트 그중에 뽈라서 실행에 올린다. 그다음에 그걸 학습해서 bake한다는 이게 그냥 옛날아이디어야 이아이디어는 무시해두되는데 이런식으로 harvest를 가져가도 좋고 다른 아이디어 있으면 그걸로 한번 해줘봐도돼 후보를 만들어보자 재치있음 좋겠어. 개인적으로 나는 이게 엄청난 혁신이 될 수도 있겠단 생각을 해서 그래}
  \item \textbf{사용자} (ul 129, 09-27 01:5x KST): \intent{PaceNotes가 제일 잘 맞는듯}
  \item \textbf{그래서(결정):} user-log 129(f04b41d): 제목 `PaceNotes: Intent-Sharing Hierarchies for Fast, Verifiable Robot Manipulation' --- 랠리 코드라이버가 앞 구간 노트를 불러 주고(상위 의도 선언), 드라이버가 실시간 조향·제동 타이밍을 잡으며(VLA 세부 조종·잡기 타이밍), 코드라이버가 위치를 확인(상위 검증)하는 비유(§96 보충 2와 같은 구조). 저장소·코드 이름 harvest는 유지.
  \item \textbf{참고:} user-log 129의 시각 표기(16:5x UTC = 01:5x KST)가 커밋 시각(01:46 KST)보다 늦다 --- 기록 그대로 둔다.
\end{itemize}

\paragraph{F10. 논문도 상위 중심, 결합은 §96 보충 5로}
{\small 09-27 03:2x KST}
\begin{itemize}
  \item \textbf{사용자} (ul 140, 03:2x KST): \intent{논문도 이런식으로 적는걸로 하자 일단 실험들은 전부 상위를 중심으로 해야하는 것은 맞구 말이야.}
  \item \textbf{그래서(결정):} 논문 da3b98f: 근거 5절을 `상위 계획기(주 줄기)'와 `2단계(조건부)'로 나눔, 3.4절·Fig. 1·모델 그림을 §96 보충 5로, 커밋 안 된 결과는 `진행 중'으로만.
\end{itemize}

\paragraph{F11. 버리기 전 재검증 + 살아난 안은 공정 비교 한 번 더}
{\small 09-27 03:4x KST}
\begin{itemize}
  \item \textbf{사용자} (ul 143, 03:4x KST): \intent{결정 버리기전에 재검증해서 다시 살아났다고 기존거 ㄹ또 버리지말고 두개 비교를 공정하게 한번더 하는과정이 있어야해 알겠지?}
  \item \textbf{그래서(결정):} 협업 보드 규칙·메모리(b02ff17): 재검증으로 살아난 안과 기존 안은 같은 데이터·학습량·평가 세트·시드·판정 규칙으로 사전 등록한 짝 비교를 한 번 더 한 뒤에만 선택(E-STRIP8 손 정보, T3 재검증, 작전 T). 논문 6절 절차 규칙에 반영.
\end{itemize}

